%% Elements of Nested Fibrational Cosmology
%% CRYST Branch Book: Crystallography Branch (Prospective)
%% Status: Architecturally constituted; pre-closure.
%%
%% Status legend:
%%   [U]  Unconditional  -- proved from spine imports and prior [U] results
%%   [C]  Conditional    -- depends on explicitly declared hypotheses
%%   [D]  Definition-structural
%%   [R]  Remark only
%%   [O]  Open obligation / named failure frontier
%%   [B]  Bridge (Book VI licensed continuum surrogate)
%%
%% Notation:
%%   R_Cryst              Crystallography target regime
%%   O_Cryst              Crystallography observable family = (Phi_Cryst, L_Cryst, G_Cryst)
%%   Phi_Cryst            Diffraction kernel (quotient-visible intensity data)
%%   L_Cryst              Lattice ledger (certified translation equivalence structure)
%%   G_Cryst              Symmetry-descent object (certified space group carrier)
%%   B_Cryst(X)           Bragg gap (positive peak separation from diffuse background)
%%   F_Cryst(hkl)         Structure factor (amplitude; largely quotient-invisible; see phase problem)
%%   E_Cryst              Crystallography endpoint datum
%%   I(hkl)               Bragg intensity observable = |F_Cryst(hkl)|^2 (quotient-visible)
%%   SG_Cryst             Certified space group object (one of 230 at declared scope)

\documentclass[12pt,oneside]{amsbook}
\usepackage{amsmath,amssymb,amsthm}
\usepackage{mathrsfs}
\usepackage{enumitem}
\usepackage{geometry}
\usepackage{hyperref}
\geometry{margin=1.2in}

\theoremstyle{plain}
\newtheorem{theorem}{Theorem}[section]
\newtheorem{lemma}[theorem]{Lemma}
\newtheorem{proposition}[theorem]{Proposition}
\newtheorem{corollary}[theorem]{Corollary}

\theoremstyle{definition}
\newtheorem{definition}[theorem]{Definition}
\newtheorem{standingrule}[theorem]{Standing Rule}
\newtheorem{openobligation}[theorem]{Open Obligation}

\theoremstyle{remark}
\newtheorem{remark}[theorem]{Remark}
\newtheorem{scholium}[theorem]{Scholium}

\newcommand{\status}[1]{\textup{[#1]}}

\begin{document}
\emergencystretch=3em
\raggedbottom

\title{NFC Crystallography Branch Book}
\author{Elements of Nested Fibrational Cosmology\\
        Prospective Derived Branch}
\date{}
\maketitle
\tableofcontents
\newpage

%% ---------------------------------------------------------------
\section{Legitimacy Analysis and Branch Constitution}
\label{sec:cryst-legitimacy}
%% ---------------------------------------------------------------

Before constituting the branch, this section records the systematic
legitimacy analysis required by Book~V's Certified Branch Legitimacy
Theorem (\texttt{thm:V.8.1}) and the UBLT (\texttt{thm:V.7.4}).
The analysis is the basis for the verdict in
Prop.~\ref{prop:cryst-lawful}; the branch is constituted only if
all five UBLT conditions are satisfied.

\subsection{The Question}

\begin{remark}[\status{R}]\label{rem:cryst-question}
\textbf{Crystallography and NFC: the central question.}
Standard crystallography operates on three imported primitives that
the NFC framework forbids: (i) Euclidean three-dimensional geometry
(a continuous background space); (ii) Fourier analysis (the
reciprocal lattice and structure factors); (iii) atomic species (Na,
Cl, C, \ldots\ as chemical primitives carrying pre-given scattering
properties).  The central question of this branch is whether
crystallographic structure can be realized as a lawful descendant of
$\mathbf{U}$ from purely finite observational primitives, without any
of these three imports.

Specifically: starting from finite probe-response pairs (a beam
probe, an intensity response), a finite admissible contrast discipline
on the response data, and the observational equivalence imposed by
that discipline, can one recover certified lattice periodicity, space
group symmetry, and diffraction peak structure?  Each of these
notions, if it appears at all, must emerge as a quotient-visible,
source-descended object within the declared regime.

The answer, given in Prop.~\ref{prop:cryst-lawful}, is: the
crystallography program constitutes a lawful branch candidate.  All
five UBLT conditions can be satisfied at the declared pre-closure
scope.
\end{remark}

\subsection{Five UBLT Conditions: Pre-Analysis}

\begin{remark}[\status{R}]\label{rem:cryst-ublt-analysis}
\textbf{Pre-analysis of Book~V UBLT conditions.}
The five UBLT conditions (Book~V, \texttt{thm:V.7.4}) and their
crystallographic realizations:

\medskip
\textbf{(1) Declared branch candidate.}
The crystallography branch candidate $\mathcal{B}_{\mathrm{Cryst}}$
is the program targeting certified diffraction-peak structure,
lattice periodicity, and space group symmetry on the declared regime
$\mathcal{R}_{\mathrm{Cryst}}$.  The five mandatory components
(realization functor $F_{\mathrm{Cryst}}$, burden package
$\mathrm{Cert}_{\mathrm{Cryst}}$, terminal predicate
$\tau_{\mathrm{Cryst}}$, failure frontier
$\mathrm{FF}_{\mathrm{Cryst}}$, and observable family
$\mathcal{O}_{\mathrm{Cryst}}$) are declared below.  \emph{Condition
(1): satisfied.}

\medskip
\textbf{(2) Observable family visibility.}
The primary diffraction observable is the Bragg intensity
$I(hkl) = |F_{\mathrm{Cryst}}(hkl)|^2$.  This is the squared
modulus of the structure factor: it is a nonnegative real number
measured at each reciprocal lattice point $(hkl)$ and is directly
observable under admissible diffraction probes.  The intensity
is quotient-visible: two crystal configurations indistinguishable
by all admissible diffraction probes have the same intensity profile.
The Bragg gap $B_{\mathrm{Cryst}} > 0$ (positive separation of
Bragg peaks from the diffuse background) is also quotient-visible:
it is the minimum intensity contrast between the peak maximum and
the inter-peak floor, measurable from the certified intensity profile.
\emph{Condition (2): satisfied.}

\medskip
\textbf{(3) Source-descended comparison.}
The diffraction probe-response structure descends from $\mathbf{U}$
via the same route as the SPEC branch: probe episodes are admissible
processes on the certified observable class; their responses are
certified observable outcomes; the comparison discipline is the
observational quotient of Book~I applied to intensity profiles.
The lattice periodicity structure descends from $\mathbf{U}$ via
the Book~III transport law applied to the admissible equivalence
under lattice translations.  \emph{Condition (3): satisfied.}

\medskip
\textbf{(4) No hidden supplement.}
Three potential hidden imports are identified and addressed:
\begin{enumerate}[label=\textup{(\alph*)}]
  \item \emph{Euclidean geometry}: not imported as primitive.  The
    three-dimensional spatial structure of the crystal must emerge
    as the effective encoding of certified discrete lattice data,
    licensed by a Book~VI bridge (O-CRYST.EWALD).
  \item \emph{Fourier analysis}: the reciprocal lattice and
    structure factor notation are licensed only after the discrete
    diffraction peaks are certified and the Book~VI asymptotic
    coherence bridge (O-CRYST.EWALD) is established.
  \item \emph{Atomic species}: atoms are not primitives.  A
    \emph{scatterer} in this branch is characterized solely by
    its certified probe-response behavior (scattering length or
    electron density contribution), not by its chemical identity.
    Chemical identity, if it appears at all, must emerge from the
    observable scattering data.
\end{enumerate}
\emph{Condition (4): satisfied (pending O-CRYST.EWALD for the
Fourier/Euclidean bridge).}

\medskip
\textbf{(5) No path multiplication.}
The crystallography branch's terminal predicate (certified Bragg
peak structure with positive gap and space group descent) is not
the terminal predicate of any existing branch:
\begin{itemize}
  \item YM targets a spectral mass gap for the covariant Laplacian
    on the gauge sector --- distinct from a diffraction gap;
  \item SPEC targets probe-response resonance structure --- adjacent
    but not identical (SPEC does not target periodicity or space
    group structure);
  \item BIO targets replication-heredity structure --- distinct;
  \item LING targets contrast-compositional structure --- distinct;
  \item SCC targets structural counterfactual capacity --- distinct.
\end{itemize}
The crystallography terminal predicate is distinct from all existing
named branches.  \emph{Condition (5): satisfied.}
\end{remark}

\subsection{Legitimacy Proposition}

\begin{proposition}[\status{U}]\label{prop:cryst-lawful}
\textup{(Crystallography Branch Is Lawful.)}
The crystallography branch candidate $\mathcal{B}_{\mathrm{Cryst}}$
satisfies all five conditions of the Certified Branch Legitimacy
Theorem (Book~V, \texttt{thm:V.8.1}).  It is a lawful NFC branch
candidate.
\end{proposition}

\begin{proof}
Each of the five UBLT conditions is addressed in
Rem.~\ref{rem:cryst-ublt-analysis}: (1) all five mandatory
components are declared in §\ref{sec:cryst-purpose}--\ref{sec:cryst-closure};
(2) the Bragg intensity $I(hkl)$ is quotient-visible and the Bragg
gap $B_{\mathrm{Cryst}} > 0$ is branch-visible;
(3) source descent follows the SPEC branch template via Book~I
observational quotient and Book~III transport law;
(4) no Euclidean geometry, Fourier analysis, or atomic species is
imported as primitive --- each is either earned by a Book~VI bridge
or dismissed by the no-smuggling standing rules;
(5) the terminal predicate is distinct from all existing branches
by the terminal-predicate distinctness argument.
By Book~V \texttt{thm:V.8.1}, $\mathcal{B}_{\mathrm{Cryst}}$ is
a lawful branch candidate.
\qed
\end{proof}

\begin{proposition}[\status{U}]\label{prop:cryst-distinct}
\textup{(Full Distinctness.)}
The crystallography branch is genuinely distinct from all existing
named branches.  Its terminal predicate --- certified Bragg peak
structure with positive gap $B_{\mathrm{Cryst}} > 0$ and
source-descended space group classification --- is not the terminal
predicate of YM, NS, GR, SCC, RH, SM, SPEC, BIO, or LING.
\end{proposition}

\begin{proof}
By the pairwise comparison in Rem.~\ref{rem:cryst-ublt-analysis}(5):
the Bragg gap is not the YM mass gap (different carrier and probe
family); not the SPEC spectral margin (SPEC targets resonance
structure, not lattice periodicity); not the BIO fidelity margin
(different observable class); not the LING contrast margin
(different branch regime).  No other branch targets space group
symmetry as its terminal predicate.
\qed
\end{proof}

\begin{proposition}[\status{U}]\label{prop:cryst-anti-smuggle}
\textup{(Crystallography Anti-Smuggling Gate.)}
No branch claim may import Euclidean geometry, Fourier analysis,
or atomic species as branch primitives.  These must be either (a)
derived as effective encodings of certified discrete observable
structure via Book~VI bridges, or (b) dismissed by the
no-smuggling standing rules.
\end{proposition}

\begin{proof}
By Book~I Standing Rule~I.1.3 (no-smuggling at the primitive level)
and the branch-level extension (Book~V UBLT condition~(4)):
any datum outside the certified source image $F_{\mathrm{Cryst}}(\mathbf{U})$
is a silent import and violates the intake law.
\qed
\end{proof}

%% ---------------------------------------------------------------
\section{Purpose and Canonical Seed}
\label{sec:cryst-purpose}
%% ---------------------------------------------------------------

This branch asks whether crystallographic structure --- lattice
periodicity, space group symmetry, diffraction peak separation, and
structure determination --- can be realized as a lawful descendant of
the NFC source architecture.  The aim is not to import the standard
language of Euclidean lattices, Fourier transforms, Miller indices,
or atomic coordinates as primitives, but to determine when such
language is licensed as an effective encoding of certified discrete
observable structure within a declared branch regime.  This posture
is mandated by Book~I's no-smuggling stance, Book~V's
branch-legitimacy law, Book~VI's continuum-interface discipline, and
Book~VII's scoped-certificate law.

The branch has three structural precedents in the existing corpus.
First, the SPEC branch already treats probe-response pairs and
transition ledgers as the building blocks of a certified observable
family --- the \emph{diffraction-probe template}: a beam probe acts
on a crystal configuration and produces an intensity response, which
is a SPEC-style probe-response pair.  Second, the BIO branch already
treats a replication-maintaining core $I_X$ and boundary-forcing
exchange as the route to a certified interface --- applicable here
because crystal grain boundaries and unit cell boundaries share the
same boundary-capacity structure.  Third, the SCC branch's minimal
structural carrier (MSC) machinery provides the anti-inflation gate
for the space group identification: among all symmetry objects
compatible with a given diffraction pattern, MSC selects the unique
minimal faithful carrier --- which, at the declared regime, is one
of the 230 crystallographic space groups.

The initial endpoint is deliberately modest: certified Bragg peak
structure on a declared crystallographic regime, with positive peak
separation from the diffuse background (the Bragg gap), and
source-descended translation equivalence classes (the lattice
structure).  Full space group identification, structure
determination, and the resolution of the phase problem are deferred
as named open obligations.

\begin{remark}[\status{R}]\label{rem:cryst-strategic-path}
\textbf{Strategic path.}
The shortest canonical route is: probe-response certification
(diffraction as admissible probe) $\Rightarrow$ Bragg gap
(positive peak separation) $\Rightarrow$ lattice certification
(translation equivalence classes) $\Rightarrow$ space group descent
(MSC-selected symmetry carrier) $\Rightarrow$ endpoint.  The phase
problem and structure determination are downstream open obligations,
expected to require Book~VI bridges and possibly new
probe-family extensions beyond intensity-only measurement.
\end{remark}

%% ---------------------------------------------------------------
\section{Imported Spine Structure}
\label{sec:cryst-imports}
%% ---------------------------------------------------------------

\textbf{Imported spine results.}
\begin{enumerate}[label=\textup{(\arabic*)}]
  \item \emph{Book~I}: observational quotient discipline; no Euclidean
    geometry, Fourier basis, atomic species, or crystallographic
    convention may be assumed.  Crystal content must be
    quotient-visible and admissibly testable.
  \item \emph{Book~II}: boundary capacity law.  Crystal grain
    boundaries and unit cell boundaries are interface structures
    whose information capacity is logarithmically bounded.  This
    governs the admissible probe resolution of crystal interfaces.
  \item \emph{Book~III}: unified coercive-transport inequality.
    Stability costs distribute across the crystal lattice without
    hidden channel.  Lattice phonons and thermal diffuse scattering
    are accounted for as defect contributions to the transport
    ledger.
  \item \emph{Book~IV}: universal source $\mathbf{U}$ and
    source-descended comparison.  All certified crystallographic
    objects factor through $\mathbf{U}$.
  \item \emph{Book~V}: lawful branchhood; the constitutional law
    under which the crystallography branch is admitted
    (Prop.~\ref{prop:cryst-lawful}).
  \item \emph{Book~VI}: continuum language licensed only as faithful
    shorthand for certified discrete structure within a declared
    interface regime.  The Book~VI gate controls:
    continuous angle/wavelength notation (Bragg's law as
    asymptotic shorthand), Fourier/reciprocal lattice notation
    (licensed only after O-CRYST.EWALD is discharged), and
    Euclidean $\mathbb{R}^3$ domain notation (licensed only after
    O-CRYST.DOM is discharged).
  \item \emph{Book~VII}: scoped certificates, promotion law, named
    frontier stability, and prohibition of hidden upgrade.  Claims
    about universal crystallographic laws, space group uniqueness,
    or structure determination from intensities alone are subject to
    this governance without exception.
\end{enumerate}

\textbf{Imported branch precedents.}
\begin{enumerate}[label=\textup{(\arabic*)}]
  \item \emph{SPEC branch}: probe-response objects and transition
    ledgers --- the \emph{diffraction-probe template} for this branch.
    Diffraction is a SPEC-style probe episode; intensity data is a
    SPEC-style response observable.
  \item \emph{BIO branch}: boundary-forcing exchange and the
    replication-maintaining core --- the \emph{unit-cell boundary
    template}.  A crystal unit cell's interior/exterior separation
    follows the same boundary-forcing logic as BIO's
    $\partial I_X$.
  \item \emph{SCC branch}: MSC machinery --- the
    \emph{space-group minimality template}.  Among all
    certified symmetry objects compatible with a given intensity
    profile, MSC selects the unique minimal faithful carrier.
  \item \emph{YM branch}: gap structure --- the Bragg gap
    $B_{\mathrm{Cryst}} > 0$ is the fifth canonical instance of
    the positive-margin separation pattern (after YM mass gap,
    SPEC spectral margin, LING contrast margin, and BIO fidelity
    margin).
\end{enumerate}

%% ---------------------------------------------------------------
\section{Branch-Specific Arena}
\label{sec:cryst-arena}
%% ---------------------------------------------------------------

\subsection{Standing Rules}

\begin{standingrule}[\status{D}]\label{sr:cryst-no-smuggling}
\textup{(No Smuggling.)}
No crystallographic object --- Bravais lattice, space group,
Miller index, structure factor, atomic coordinate, or unit cell ---
may be treated as primitive branch data merely because it is
familiar in external crystallographic theory.  Every branch-visible
crystallographic object must be sourced through $\mathbf{U}$,
the declared descent machinery, and the crystallography observable
family.  This is the branch-local instance of Book~I no-smuggling
and Book~V branch-legitimacy discipline.
\end{standingrule}

\begin{standingrule}[\status{D}]\label{sr:cryst-no-euclid}
\textup{(No Euclidean Primitive.)}
Euclidean three-dimensional geometry is not a primitive of this
branch.  The ambient space of the crystal, if it appears at all,
must emerge as the effective encoding of certified discrete lattice
data via a Book~VI domain bridge (O-CRYST.DOM).  Until that bridge
is established, all spatial language is forbidden as a branch
primitive.
\end{standingrule}

\begin{standingrule}[\status{D}]\label{sr:cryst-no-fourier}
\textup{(No Fourier Primitive.)}
Fourier analysis, the reciprocal lattice, structure factors as
complex amplitudes, and Miller index notation as geometric primitive
are not licensed until the Book~VI Ewald bridge (O-CRYST.EWALD)
is established.  The only immediately certified diffraction
observable is the Bragg intensity $I(hkl) = |F(hkl)|^2$, which is
a real nonnegative probe-response datum --- not the full complex
structure factor.
\end{standingrule}

\begin{standingrule}[\status{D}]\label{sr:cryst-no-atom}
\textup{(No Atomic Primitive.)}
No atomic species, chemical element, or molecular formula may
be treated as a primitive.  A \emph{scatterer} in this branch is
characterized solely by its certified probe-response behavior within
the declared admissible probe discipline.  Chemical identity, if it
appears at all, must emerge from the observable scattering data.
\end{standingrule}

\begin{standingrule}[\status{D}]\label{sr:cryst-quotient-visible}
\textup{(Quotient Visibility.)}
A crystallographic distinction between two configurations counts
only when it is quotient-visible under admissible diffraction probe
tests.  Two crystal configurations whose intensity profiles agree
on all admissible diffraction probes belong to the same certified
crystallographic equivalence class, regardless of their chemical
or geometric description in external notation.
\end{standingrule}

\begin{standingrule}[\status{D}]\label{sr:cryst-no-self-promotion}
\textup{(No Self-Promotion.)}
A finite certified diffraction result may not self-promote to a
claim about all crystal structures, universal space group
uniqueness, or complete structure determination from intensities
alone.  Such promotion requires an explicit transfer theorem.
This is a direct Book~VII governance consequence.
\end{standingrule}

\begin{standingrule}[\status{D}]\label{sr:cryst-frontier-visibility}
\textup{(Frontier Visibility.)}
Named frontiers in space group identification, structure
determination, the phase problem, domain certification, and
Fourier legitimacy remain visible until discharged by theorem.
Rephrasing does not discharge them.
\end{standingrule}

\begin{standingrule}[\status{D}]\label{sr:cryst-continuum-gate}
\textup{(Continuum Gate.)}
Continuous angle notation ($2\theta$), wavelength notation
($\lambda$), Fourier transform notation for electron density, and
Euclidean coordinate notation may appear only after the
corresponding Book~VI bridge theorem certifies that they faithfully
encode certified discrete observable structure within the declared
continuum interface regime.
\end{standingrule}

\subsection{Declared Target Regime}

\begin{definition}[\status{D}]\label{def:cryst-target-regime}
\textup{(Crystallography Target Regime.)}
A \emph{crystallography target regime} is a declared branch regime
$\mathcal{R}_{\mathrm{Cryst}}$ consisting of:
\begin{enumerate}[label=\textup{(\arabic*)}]
  \item a certified observable class $\mathcal{C}_{\mathrm{obs}}$
    of physical configurations --- the objects on which diffraction
    probe tests are performed (crystal specimens in the declared
    regime, not atoms or molecules);
  \item a finite admissible diffraction probe discipline
    $\mathcal{P}_{\mathrm{diff}}$, each element of which is a
    declared admissible process producing a
    quotient-visible intensity response;
  \item a declared probe-energy window $\mathcal{W}_{\mathrm{diff}}$
    specifying the range of probe energies (wavelengths) within which
    diffraction claims have force; and
  \item a declared comparison discipline governing when two
    intensity profiles $\{I(hkl)\}$ are equivalent for the purpose
    of crystallographic claims.
\end{enumerate}
\end{definition}

\begin{remark}[\status{R}]\label{rem:cryst-configuration-note}
\textbf{What a configuration is.}
A \emph{physical configuration} in this branch is a certified
observable arrangement of scatterers (not atoms or molecules ---
those are chemical primitives) on which admissible diffraction
probes can be performed.  Whether any such configuration qualifies
as periodic, crystalline, or amorphous is determined by the
diffraction kernel and lattice ledger, not by prior crystallographic
classification.
\end{remark}

%% ---------------------------------------------------------------
\section{Crystallography Observable Family}
\label{sec:cryst-obs-family}
%% ---------------------------------------------------------------

\subsection{Three-Part Structure}

\begin{definition}[\status{D}]\label{def:cryst-obs-family}
\textup{(Crystallography Observable Family.)}
A \emph{crystallography observable family} on
$\mathcal{R}_{\mathrm{Cryst}}$ is a triple
\[
  \mathcal{O}_{\mathrm{Cryst}} \;=\;
  \bigl(\Phi_{\mathrm{Cryst}},\;
        \mathcal{L}_{\mathrm{Cryst}},\;
        \mathcal{G}_{\mathrm{Cryst}}\bigr)
\]
where:
\begin{enumerate}[label=\textup{(\arabic*)}]
  \item $\Phi_{\mathrm{Cryst}}$ is the \emph{diffraction kernel}: the
    quotient-visible intensity profile induced on
    $\mathcal{C}_{\mathrm{obs}}$ by the admissible diffraction probe
    discipline $\mathcal{P}_{\mathrm{diff}}$;
  \item $\mathcal{L}_{\mathrm{Cryst}}$ is the \emph{lattice ledger}:
    certified records of translation equivalence classes on the
    observable configuration, with full transport bookkeeping; and
  \item $\mathcal{G}_{\mathrm{Cryst}}$ is the \emph{symmetry-descent
    object}: a branch-visible certified symmetry carrier governing
    the complete set of admissible symmetry operations on the
    diffraction-visible configuration.
\end{enumerate}
\end{definition}

\subsection{Diffraction Kernel}

\begin{definition}[\status{D}]\label{def:cryst-diffraction-kernel}
\textup{(Diffraction Kernel.)}
The \emph{diffraction kernel} $\Phi_{\mathrm{Cryst}}$ is the
assignment
\[
  \Phi_{\mathrm{Cryst}} : (X, P_{hkl}) \;\mapsto\;
  I(hkl) := |F(hkl)|^2
\]
from a configuration $X$ and admissible diffraction probe
$P_{hkl} \in \mathcal{P}_{\mathrm{diff}}$ to the certified
intensity response.  Two configurations belong to the same
\emph{diffraction equivalence class} $[X]_{\Phi}$ if and only if
no admissible probe in $\mathcal{P}_{\mathrm{diff}}$ distinguishes
their intensity profiles.
\end{definition}

\begin{remark}[\status{R}]\label{rem:cryst-phase-note}
\textbf{The phase is not part of the kernel.}
The diffraction kernel records $I(hkl) = |F(hkl)|^2$ --- the
squared modulus of the structure factor.  The phase
$\phi(hkl) = \arg F(hkl)$ is not part of the intensity-observable
kernel.  Whether phase information is quotient-invisible (cannot
be recovered from intensity data alone under any admissible probe
discipline) or recoverable via richer probes is the content of the
principal frontier O-CRYST.PHASE (§\ref{sec:cryst-phase}).  This
is the most structurally interesting feature of the crystallography
branch: the \emph{phase problem} maps directly onto the NFC
question of what is and is not quotient-visible under the declared
admissible probe discipline.
\end{remark}

\begin{definition}[\status{D}]\label{def:cryst-bragg-gap}
\textup{(Bragg Gap.)}
The \emph{Bragg gap} at a configuration $X$ is
\[
  B_{\mathrm{Cryst}}(X) \;:=\;
  \min_{hkl \in \mathrm{peak}(X)}\;
  \bigl[\,I(hkl) - I_{\mathrm{bg}}(X)\,\bigr],
\]
where $\mathrm{peak}(X)$ is the certified set of Bragg peak
indices and $I_{\mathrm{bg}}(X)$ is the certified diffuse
background intensity at those positions.  The Bragg gap is
positive when diffraction peaks are discretely separated from the
background: $B_{\mathrm{Cryst}}(X) > 0$.
\end{definition}

\begin{remark}[\status{R}]\label{rem:cryst-gap-analogy}
\textbf{The Bragg gap as the fifth canonical gap instance.}
$B_{\mathrm{Cryst}}(X) > 0$ is the crystallography branch's instance
of the canonical positive-margin separation structure that appears
across the NFC corpus: the YM mass gap ($m^2 > 0$), the SPEC
spectral margin ($\lambda_1 > 0$), the LING contrast margin
($\Delta_{\mathrm{Ling}} > 0$), and the BIO replication fidelity
margin ($F_{\mathrm{Bio}} > 0$).  In each case the structural claim
is a positive observable separation between distinct certified
equivalence classes.  Here: Bragg peaks are discretely separated
from the diffuse scattering background in a certified
intensity-observable sense.
\end{remark}

\subsection{Lattice Ledger}

\begin{definition}[\status{D}]\label{def:cryst-lattice-ledger}
\textup{(Lattice Ledger.)}
The \emph{lattice ledger} $\mathcal{L}_{\mathrm{Cryst}}$ is a
certified bookkeeping object recording, for a configuration $X$:
\begin{enumerate}[label=\textup{(\arabic*)}]
  \item the certified set of translation equivalence classes
    $\{[x]_{\sim_t}\}$ --- the quotient of the configuration's
    admissible sites under the certified translation probe discipline;
  \item the certified transport cost and defect contribution of each
    admissible translation, governed by Book~III; and
  \item the Book~II boundary capacity bookkeeping for the unit cell
    boundary $\partial U_X$ (the interface between one translation
    equivalence class and its neighbors).
\end{enumerate}
The lattice ledger is the crystallographic analogue of the
compositional ledger (LING branch) and the metabolic ledger (BIO
branch).
\end{definition}

\begin{remark}[\status{R}]\label{rem:cryst-translation-equiv}
\textbf{Lattice as quotient structure.}
The Bravais lattice, in standard crystallographic terms, is the
group of translation vectors that map the crystal onto itself.
In NFC terms, it is the quotient of the configuration's admissible
sites under the certified translation probe equivalence: two sites
$x, x'$ are translation-equivalent if no admissible diffraction
probe distinguishes them when one is mapped to the other by an
admissible process.  The lattice is derived, not assumed.  Whether
it is one of the 14 Bravais lattice types is a derived theorem
requiring Book~VI geometric legitimacy, not a branch assumption.
\end{remark}

\subsection{Symmetry-Descent Object}

\begin{definition}[\status{D}]\label{def:cryst-symmetry-descent}
\textup{(Symmetry-Descent Object.)}
A \emph{symmetry-descent object} $\mathcal{G}_{\mathrm{Cryst}}$ is
a branch-visible certified symmetry carrier that:
\begin{enumerate}[label=\textup{(\arabic*)}]
  \item is sourced through $\mathbf{U}$;
  \item is selected by MSC/minimal-support discipline as the unique
    minimal faithful carrier for the diffraction-visible symmetry
    predicate on the declared regime;
  \item produces only quotient-visible distinctions among
    configurations; and
  \item passes the Book~V five-condition branch-legitimacy test.
\end{enumerate}
$\mathcal{G}_{\mathrm{Cryst}}$ is not a space group declared by
external crystallographic convention.  It is a derived certified
carrier.  If, under the declared regime and admissible probe
discipline, MSC selects a carrier matching one of the 230
crystallographic space groups, that match is a derived theorem ---
the content of O-CRYST.SG --- not a branch assumption.
\end{definition}

\begin{remark}[\status{R}]\label{rem:cryst-230-note}
\textbf{The 230 space groups as a prediction, not an assumption.}
Standard crystallography enumerates 230 crystallographic space
groups as the complete set of symmetry groups compatible with
three-dimensional Euclidean lattice periodicity.  In the NFC
framework, this enumeration must be derived: the claim is that
MSC applied to the diffraction-visible symmetry predicate on the
declared Euclidean domain (after O-CRYST.DOM and O-CRYST.EWALD
are discharged) selects carriers that, up to branch-visible
equivalence, coincide with one of these 230 groups.  This is a
precise, falsifiable prediction of the branch --- its compact
derivable content.  If MSC selects more or fewer carriers than
230, the branch has discovered either a new symmetry class or an
error in the existing enumeration.
\end{remark}

%% ---------------------------------------------------------------
\section{Endpoint Datum and Branch Posture}
\label{sec:cryst-endpoint}
%% ---------------------------------------------------------------

\begin{definition}[\status{D}]\label{def:cryst-endpoint}
\textup{(Crystallography Endpoint Datum.)}
The \emph{crystallography endpoint datum} is the branch-visible
package
\[
  E_{\mathrm{Cryst}} \;=\;
  \bigl(\mathcal{C}_{\mathrm{obs}},\;
        \mathcal{O}_{\mathrm{Cryst}},\;
        B_{\mathrm{Cryst}},\;
        \mathcal{W}_{\mathrm{diff}}\bigr)
\]
consisting of the certified observable class, the crystallography
observable family, the Bragg gap, and the declared probe-energy
window.
\end{definition}

\begin{definition}[\status{D}]\label{def:cryst-lawful-descendant}
\textup{(Lawful Crystallography Descendant Claim.)}
A \emph{lawful crystallography descendant claim} is a claim that
the endpoint datum $E_{\mathrm{Cryst}}$ is sourced through
$\mathbf{U}$ and supports a certified diffraction, periodicity,
or symmetry statement on the declared regime.
\end{definition}

\begin{proposition}[\status{C}]\label{prop:cryst-branch-candidate}
\textup{(Crystallography Branch Candidate.)}
\textup{[Conditional on the crystallography target regime, observable
family, and endpoint datum all being sourced through $\mathbf{U}$,
branch-visible, and admitting a Book~V legitimacy witness.]}
The crystallography program is a lawful branch candidate.
\end{proposition}

\begin{proof}
By Prop.~\ref{prop:cryst-lawful} and the branch-legitimacy doctrine
(Book~V) applied to the declared crystallography endpoint datum.
\qed
\end{proof}

%% ---------------------------------------------------------------
\section{Core Theorem Shells and Named Frontiers}
\label{sec:cryst-theorems}
%% ---------------------------------------------------------------

\subsection{Probe Realization}

\begin{openobligation}[\status{C}]\label{ob:cryst-PR}
\textup{(O-CRYST.PR: Probe Realization.)}
Show that the diffraction probe family $\mathcal{P}_{\mathrm{diff}}$
is sourced through $\mathbf{U}$ and constitutes a lawful branch
observable within the declared regime, with all probe declarations
inherited from the Book~I admissible-test discipline.
\end{openobligation}

\begin{theorem}[\status{C}]\label{thm:cryst-PR}
\textup{(Probe Realization Theorem.)}
\textup{[Conditional on O-CRYST.PR.]}
The diffraction kernel $\Phi_{\mathrm{Cryst}}$ is a lawful branch
observable descended from $\mathbf{U}$.  The certified diffraction
equivalence classes $\{[X]_{\Phi}\}$ are quotient-visible and
admissibly testable.
\end{theorem}

\begin{proof}
By the branch-legitimacy doctrine (Book~V) applied to the declared
diffraction kernel, together with the Book~I observational quotient
discipline.  Full proof is conditional on O-CRYST.PR.
\qed
\end{proof}

\subsection{Bragg Gap}

\begin{openobligation}[\status{C}]\label{ob:cryst-GAP}
\textup{(O-CRYST.GAP: Bragg Gap.)}
Show that the Bragg gap $B_{\mathrm{Cryst}}(X)$ is
branch-computably positive for certified crystalline configurations
within the declared regime, and that no Bragg peak collapse occurs
without a declared admissible structure-change event.
\end{openobligation}

\begin{theorem}[\status{C}]\label{thm:cryst-GAP}
\textup{(Bragg Gap Theorem.)}
\textup{[Conditional on O-CRYST.PR and O-CRYST.GAP; the admissible
diffraction probe family being non-trivially separating on the
declared regime.]}
For any certified crystalline configuration $X$,
$B_{\mathrm{Cryst}}(X) > 0$: Bragg peaks are discretely separated
from the diffuse scattering background.
\end{theorem}

\begin{proof}
By the observational quotient discipline (Book~I) and the declared
admissible-test finiteness condition on $\mathcal{R}_{\mathrm{Cryst}}$.
The positivity of the Bragg gap reflects the separating power of the
admissible diffraction probe discipline.  Full proof is conditional
on O-CRYST.GAP.
\qed
\end{proof}

\subsection{Lattice Certification}

\begin{openobligation}[\status{C}]\label{ob:cryst-LAT}
\textup{(O-CRYST.LAT: Lattice Certification.)}
Show that the lattice ledger $\mathcal{L}_{\mathrm{Cryst}}$ is
additive, sourced through the Book~III transport accounting, and
that the translation equivalence classes are genuinely
quotient-visible (no hidden spatial primitive is required).
\end{openobligation}

\begin{theorem}[\status{C}]\label{thm:cryst-LAT}
\textup{(Lattice Certification Theorem.)}
\textup{[Conditional on O-CRYST.PR and O-CRYST.LAT.]}
The lattice ledger $\mathcal{L}_{\mathrm{Cryst}}$ is additive over
declared translation steps.  Transport costs distribute without
hidden channel; translation equivalence classes are
quotient-visible.
\end{theorem}

\begin{proof}
By the Book~III unified coercive-transport inequality applied to the
translation operations, together with the no-smuggling standing
rule (Sr.~\ref{sr:cryst-no-euclid}).  Full proof is conditional on
O-CRYST.LAT.
\qed
\end{proof}

\subsection{No Hidden Diffraction Channel}

\begin{openobligation}[\status{C}]\label{ob:cryst-NC}
\textup{(O-CRYST.NC: No Hidden Diffraction Channel.)}
Show that no diffraction pathway exists outside the declared
admissible probe family: any discrepancy between two configurations
sharing the same certified diffraction equivalence class would
require an undeclared diffraction channel, ruled out by the
exhaustive source-image condition.
\end{openobligation}

\begin{theorem}[\status{C}]\label{thm:cryst-NC}
\textup{(No-Hidden-Diffraction-Channel Theorem.)}
\textup{[Conditional on O-CRYST.LAT and the exhaustive source-image
condition (\texttt{def:exhaustive-source-image}, Book~IV).]}
Two configurations sharing the same certified diffraction equivalence
class have the same Bragg intensity profile under all admissible
probes.  No undeclared diffraction channel can produce a
branch-visible intensity distinction invisible to the declared
intensity ledger.
\end{theorem}

\begin{proof}
Any discrepancy would require an undeclared diffraction channel.
By the exhaustive source-image condition applied to $\mathbf{U}$
and the no-Fourier-primitive standing rule
(Sr.~\ref{sr:cryst-no-fourier}), no such channel is admitted.
Full proof is conditional on O-CRYST.NC.
\qed
\end{proof}

\subsection{Space Group Descent}

\begin{openobligation}[\status{C}]\label{ob:cryst-SG}
\textup{(O-CRYST.SG: Space Group Descent.)}
Show that the symmetry-descent object $\mathcal{G}_{\mathrm{Cryst}}$
is sourced through $\mathbf{U}$, passes the Book~V five-condition
legitimacy test, and is selected by MSC/minimal-support discipline
as the unique minimal faithful carrier for the diffraction-visible
symmetry predicate.  Under the declared Euclidean domain (after
O-CRYST.DOM and O-CRYST.EWALD), show that $\mathcal{G}_{\mathrm{Cryst}}$
coincides with one of the 230 crystallographic space groups.
\end{openobligation}

\begin{theorem}[\status{C}]\label{thm:cryst-SG}
\textup{(Space Group Descent Theorem.)}
\textup{[Conditional on O-CRYST.PR, O-CRYST.NC, O-CRYST.SG, and
the declared Euclidean domain (O-CRYST.DOM).]}
The symmetry-descent object $\mathcal{G}_{\mathrm{Cryst}}$ is a
lawful branch-visible certified symmetry carrier.  Its MSC-selected
faithful subcarrier coincides with one of the 230 crystallographic
space groups within the declared domain scope.
\end{theorem}

\begin{proof}
By the Book~V branch-legitimacy doctrine applied to
$\mathcal{G}_{\mathrm{Cryst}}$, together with the no-smuggling
and no-Euclidean-primitive standing rules.  MSC selects the minimal
faithful carrier by the Zorn/finite-stabilization argument
(\texttt{thm:YM-MSC}, YM Branch §5A, transplanted to the
crystallographic symmetry-carrier family).  The coincidence with
one of the 230 space groups requires O-CRYST.DOM (for the domain)
and O-CRYST.EWALD (for the geometric structure).  Full proof is
conditional on O-CRYST.SG.
\qed
\end{proof}

\subsection{Domain Bridge}

\begin{openobligation}[\status{C}]\label{ob:cryst-DOM}
\textup{(O-CRYST.DOM: Euclidean Domain Bridge.)}
Show that the crystallography branch has a certified
three-dimensional effective geometric carrier $D_{\mathrm{Cryst}}$,
descended from $\mathbf{U}$ via a Book~VI domain bridge, satisfying
$\dim D_{\mathrm{Cryst}} = 3$ and matching the Euclidean continuum
domain at observable/operator/no-loss quotient scope.  This is the
crystallographic analogue of the YM B2 domain bridge
(\texttt{prop:ClaySpec-DOM}, YM Branch).
\end{openobligation}

\begin{remark}[\status{R}]\label{rem:cryst-dom-note}
\textbf{Why the domain bridge is required.}
The crystal's ambient three-dimensional space is not a primitive.
It must emerge as the effective encoding of certified discrete
lattice data via Book~VI.  The domain bridge proceeds exactly as in
the YM B2 analysis: build a domain reconstruction algebra
$A^{\infty}_{\mathrm{Cryst,dom}}$ from finite certified data
(diffraction observables, lattice translation records, boundary
collar data), show that its character space contains a nontrivial
three-dimensional carrier, verify Van Hove/collar exhaustion and
topological non-load-bearing cleanup, and establish no-loss passage
from the certified discrete lattice to the Euclidean continuum.
\end{remark}

\subsection{Ewald Bridge (Fourier Legitimacy)}

\begin{openobligation}[\status{C}]\label{ob:cryst-EWALD}
\textup{(O-CRYST.EWALD: Ewald/Fourier Bridge.)}
Show that reciprocal lattice notation, structure factor notation
($F(hkl)$ as a complex amplitude), and the Bragg equation
($2d\sin\theta = n\lambda$) are licensed as effective shorthands
for certified discrete diffraction peak data within a declared
Fourier-interface regime.  This is the Book~VI bridge controlling
all continuous-angle, continuous-wavelength, and complex-amplitude
notation in the crystallography branch.
\end{openobligation}

\begin{remark}[\status{R}]\label{rem:cryst-ewald-note}
\textbf{The Ewald bridge as Book~VI instance.}
The Ewald construction in standard crystallography geometrically
relates the direct lattice, the reciprocal lattice, and the
sphere of reflection in Euclidean space.  In NFC terms, the
Ewald bridge is a Book~VI continuum interface theorem: it
licenses the reciprocal lattice as the asymptotic limit of a
certified sequence of finite diffraction peak records, and it
licenses the structure factor $F(hkl)$ as the asymptotic Fourier
coefficient of the certified scattering density.  The missing
ingredient for the Book~VI bridge is the asymptotic coherence
of the finite diffraction data in the declared Fourier-interface
regime --- the crystallographic analogue of
O-LING.REC-Bridge-$\varepsilon$ and O-BIO.EVO-Bridge-$\varepsilon$.
\end{remark}

\subsection{The Phase Problem: Principal Frontier}\label{sec:cryst-phase}



\begin{remark}[\status{R}]\label{rem:cryst-phase-analysis}
\textbf{Phase problem as an NFC structural question.}
The phase problem maps with unusual precision onto the NFC framework:

\medskip
\textbf{Why the phase problem is the principal frontier.}
Standard crystallography has developed numerous methods for
recovering phase information (Patterson methods, direct methods,
isomorphous replacement, anomalous dispersion, molecular replacement).
These are all methods for extending the admissible probe family
beyond intensity-only measurement: they add prior information
(molecular geometry, heavy atoms, multiple wavelengths).  In NFC
terms, each such method is a candidate for $\mathcal{P}'_{\mathrm{diff}}$:
a richer probe family that makes more crystallographic information
quotient-visible.

\textbf{Route (a): phase invisibility.}  If the phase is
quotient-invisible under intensity-only measurement
($\mathcal{P}_{\mathrm{diff}}$), then no intensity-only experiment
can distinguish two configurations that are related by a
Babinet-complement or twin operation.  This is a certified
structural fact about the admissible probe discipline --- a
no-hidden-channel result analogous to O-LING.NC and O-BIO.NC.

\textbf{Route (b): phase recovery.}  If the phase is recoverable
via a richer probe family $\mathcal{P}'_{\mathrm{diff}}$
(e.g., adding anomalous dispersion or multiple isomorphous
replacement), the branch must certify that
$\mathcal{P}'_{\mathrm{diff}}$ is lawful (sourced through
$\mathbf{U}$, no hidden supplementation), and that the phase
information extracted is quotient-visible under that richer family.

\textbf{The NFC contribution.}  The NFC framework makes the phase
problem precise and structural: it is not merely a practical
measurement challenge, but a question about the information capacity
of the admissible probe discipline.  The Book~II boundary capacity
law bounds the maximum information extractable from any finite
probe episode; this bounds how much phase information can ever
be certified from a finite diffraction experiment.

The phase problem is therefore the crystallography branch's
counterpart to the NS Stage-3 Clay frontier: both are genuinely
open mathematical questions that the NFC framework makes precise.
\end{remark}

\subsection{Crystallography Endpoint}

\begin{openobligation}[\status{C}]\label{ob:cryst-END}
\textup{(O-CRYST.END: Crystallography Endpoint.)}
Show that the endpoint datum $E_{\mathrm{Cryst}}$ is fully certified:
probe realization, Bragg gap, lattice certification,
no-hidden-diffraction-channel, and space group descent all
discharged.  The endpoint is reached when the crystallography branch
satisfies the shared-endgame schema of Book~VII.
\end{openobligation}

\begin{theorem}[\status{C}]\label{thm:cryst-END}
\textup{(Crystallography Endpoint Theorem.)}
\textup{[Conditional on O-CRYST.PR, O-CRYST.GAP, O-CRYST.LAT,
O-CRYST.NC, O-CRYST.SG, O-CRYST.DOM, O-CRYST.EWALD.]}
The endpoint datum
\[
  E_{\mathrm{Cryst}} \;=\;
  \bigl(\mathcal{C}_{\mathrm{obs}},\;
        \mathcal{O}_{\mathrm{Cryst}},\;
        B_{\mathrm{Cryst}},\;
        \mathcal{W}_{\mathrm{diff}}\bigr)
\]
is a certified lawful crystallography endpoint: certified
diffraction kernel, additive lattice ledger, positive Bragg gap,
no-hidden-diffraction-channel control, and space group descent
sourced through $\mathbf{U}$ and MSC-selected at the declared
Euclidean domain scope.
\end{theorem}

\begin{proof}
Combine Thms.~\ref{thm:cryst-PR}, \ref{thm:cryst-GAP},
\ref{thm:cryst-LAT}, \ref{thm:cryst-NC}, and \ref{thm:cryst-SG}.
Apply the shared-endgame schema (Book~VII, Prop.~6.11): part~(i)
is the source-controlled diffraction kernel descended from
$\mathbf{U}$; part~(ii) is the lattice ledger; part~(iii) is the
no-hidden-channel theorem together with the Bragg gap.  Full proof
is conditional on O-CRYST.PR through O-CRYST.EWALD.
\qed
\end{proof}

%% ---------------------------------------------------------------
\section{Closure Ledger}
\label{sec:cryst-closure}
%% ---------------------------------------------------------------

\begin{remark}[\status{R}]\label{rem:cryst-closure}
\textbf{Obligation and frontier summary.}

\medskip
\noindent
\begin{tabular}{|l|l|p{7.8cm}|}
\hline
\textbf{Label} & \textbf{Status} & \textbf{Description} \\
\hline
O-CRYST.PR    & [O] & Probe Realization: $\Phi_{\mathrm{Cryst}}$ sourced through $\mathbf{U}$ \\
O-CRYST.GAP   & [O] & Bragg Gap: $B_{\mathrm{Cryst}} > 0$ for crystalline configurations \\
O-CRYST.LAT   & [O] & Lattice Certification: additive ledger, no hidden spatial primitive \\
O-CRYST.NC    & [O] & No Hidden Diffraction Channel: exhaustive source-image control \\
O-CRYST.SG    & [O] & Space Group Descent: MSC-selected symmetry carrier; 230 prediction \\
O-CRYST.DOM   & [O] & Euclidean Domain Bridge: $D_{\mathrm{Cryst}} \simeq \mathbb{R}^3$ at obs/op scope \\
O-CRYST.EWALD & [O] & Ewald/Fourier Bridge: reciprocal lattice and structure factor notation \\
O-CRYST.PHASE & [O] & Phase Problem (principal frontier): phase visibility or recovery route \\
O-CRYST.END   & [O] & Endpoint: full certification of $E_{\mathrm{Cryst}}$ \\
\hline
\end{tabular}

\medskip
\textbf{Discharge order.}
The natural discharge order is:
$\mathrm{PR} \to \mathrm{GAP} \to \mathrm{LAT} \to \mathrm{NC}
\to \mathrm{DOM} \to \mathrm{EWALD} \to \mathrm{SG} \to \mathrm{END}$.
The phase problem (O-CRYST.PHASE) is a named frontier orthogonal to
the main chain: it does not gate the endpoint but represents the
deepest structural question of the branch.  END can be reached
without resolving O-CRYST.PHASE: the diffraction-periodicity-symmetry
endpoint is honest even without full structure determination.

\medskip
\textbf{Two-status reading.}
The crystallography branch admits two natural closure readings:
\begin{enumerate}[label=\textup{(\alph*)}]
  \item \emph{Diffraction-periodicity-symmetry closure} (PR--END
    discharged, PHASE open): certified Bragg peak structure,
    additive lattice ledger, Bragg gap positive, and MSC-selected
    space group carrier --- the minimal honest crystallographic
    endpoint.
  \item \emph{Structure determination closure} (PR--PHASE
    discharged): adds either certified phase invisibility (no
    intensity-only experiment determines the structure) or a
    certified richer probe family recovering phase information ---
    allowing the certified intensity and phase data together to
    determine the diffraction equivalence class uniquely.
\end{enumerate}
Posture~(a) is the correct near-term target.  Posture~(b) is the
genuine long-term ambition.

\medskip
\textbf{Cross-branch relations.}
\begin{enumerate}[label=\textup{(\arabic*)}]
  \item O-CRYST.PR imports the SPEC probe-realization template
    but has no prerequisite obligations in the SPEC branch.
  \item O-CRYST.SG imports the SCC MSC machinery but requires
    only O-SCC.MCS [C] (conditionally discharged, Session~XLII),
    not the full SCC CERT-CLOSE chain.
  \item O-CRYST.DOM imports the YM B2 domain bridge pattern
    (\texttt{prop:ClaySpec-DOM}); no YM obligation is a
    prerequisite, only the structural template.
  \item O-CRYST.EWALD is analogous to O-LING.REC-Bridge-$\varepsilon$
    (asymptotic coherence condition); both are Book~VI Fourier/
    continuum bridges.
  \item O-CRYST.PHASE has no analogue in existing branches ---
    it is the first branch in the NFC corpus where the principal
    frontier is an \emph{observational information question} (what
    can be certified from intensity-only probe data?) rather than
    a structural construction question.
\end{enumerate}
\end{remark}

%% ---------------------------------------------------------------
\section{Obligation Discharge Pass}
\label{sec:cryst-discharges}
%% ---------------------------------------------------------------

This section works through the discharge order
$\mathrm{PR} \to \mathrm{GAP} \to \mathrm{LAT} \to \mathrm{NC}
\to \mathrm{DOM} \to \mathrm{EWALD} \to \mathrm{SG} \to \mathrm{END}$,
applying NFC's own machinery wherever available.  Each discharge is
conditional; no unconditional closure is claimed.  The phase problem
(O-CRYST.PHASE) receives its own formal packet in
§\ref{sec:cryst-phase-packet}.

\subsection{O-CRYST.PR: Probe Realization}
\label{sec:cryst-PR-discharge}

\begin{definition}[\status{D}]\label{def:cryst-probe-episode}
\textup{(Crystallographic Probe Episode.)}
A \emph{crystallographic probe episode} is an admissible process
$\pi = (X, P_{hkl}, R)$ consisting of a certified observable
configuration $X \in \mathcal{C}_{\mathrm{obs}}$, an admissible
probe $P_{hkl} \in \mathcal{P}_{\mathrm{diff}}$, and a
quotient-visible intensity response $R = I(hkl) \in [0,\infty)$.
Each component descends from $\mathbf{U}$; none is imported as a
pre-given quantity.
\end{definition}

\begin{definition}[\status{D}]\label{def:cryst-probe-response-object}
\textup{(Crystallographic Probe-Response Object.)}
A \emph{crystallographic probe-response object} is a SPEC-style
pair $(P_{hkl}, I(hkl))$ --- an admissible probe from
$\mathcal{P}_{\mathrm{diff}}$ together with its certified intensity
response --- viewed as a branch-visible observable datum.  The
probe-response object is lawful when both components descend from
$\mathbf{U}$ and neither introduces hidden data beyond the declared
probe discipline.
\end{definition}

\begin{theorem}[\status{C}]\label{thm:cryst-PR-discharged}
\textup{(O-CRYST.PR Conditionally Discharged.)}
\textup{[Conditional on: the admissible diffraction probe discipline
$\mathcal{P}_{\mathrm{diff}}$ being sourced through $\mathbf{U}$
(Book~IV); the Book~I observational quotient applied to intensity
profiles; and the SPEC branch probe-realization template
(\texttt{thm:spec-PR-discharged}, SPEC Branch).]}
The diffraction kernel $\Phi_{\mathrm{Cryst}}$ is a lawful branch
observable descended from $\mathbf{U}$.  The certified diffraction
equivalence classes $\{[X]_{\Phi}\}$ are quotient-visible and
admissibly testable.  O-CRYST.PR is conditionally discharged.
\end{theorem}

\begin{proof}
The proof follows the SPEC branch probe-realization template
(\texttt{thm:spec-PR-discharged}).  Each diffraction probe episode
$\pi = (X, P_{hkl}, I(hkl))$ is an admissible process on the
certified observable class.  The probe $P_{hkl}$ is a declared
admissible test (Book~I, Def.~I.1.1) with finite observable output
$I(hkl) \in [0,\infty)$.  The intensity response is quotient-visible
by the Book~I observational equivalence discipline: two configurations
producing identical intensity profiles under all probes in
$\mathcal{P}_{\mathrm{diff}}$ are observationally equivalent.  The
diffraction kernel $\Phi_{\mathrm{Cryst}}$ is the resulting
quotient-visible observable function.  No Fourier analysis, reciprocal
lattice structure, or complex amplitude is imported; only the
real-valued intensity is certified at this stage.
\qed
\end{proof}

\subsection{O-CRYST.GAP: Bragg Gap}
\label{sec:cryst-GAP-discharge}

\begin{definition}[\status{D}]\label{def:cryst-diffuse-background}
\textup{(Certified Diffuse Background.)}
The \emph{certified diffuse background} $I_{\mathrm{bg}}(X)$ at
configuration $X$ is the certified quotient-visible intensity floor:
the infimum of intensity values at probe positions not belonging to
the certified peak set $\mathrm{peak}(X)$.  It is branch-visible and
admits no hidden channel contribution not declared in the intensity
ledger.
\end{definition}

\begin{definition}[\status{D}]\label{def:cryst-periodic-config}
\textup{(Certified Crystalline Configuration.)}
A configuration $X \in \mathcal{C}_{\mathrm{obs}}$ is
\emph{crystalline} within the declared regime when its intensity
profile $\{I(hkl)\}$ under $\mathcal{P}_{\mathrm{diff}}$ has at
least one certified peak set $\mathrm{peak}(X) \neq \emptyset$ with
$B_{\mathrm{Cryst}}(X) > 0$.  This is a quotient-visible criterion:
no crystalline classification is assumed; it is derived from the
intensity data.
\end{definition}

\begin{theorem}[\status{C}]\label{thm:cryst-GAP-discharged}
\textup{(O-CRYST.GAP Conditionally Discharged.)}
\textup{[Conditional on O-CRYST.PR discharged
(Thm.~\ref{thm:cryst-PR-discharged}); the declared regime
$\mathcal{R}_{\mathrm{Cryst}}$ being non-trivially separating (at
least one probe in $\mathcal{P}_{\mathrm{diff}}$ produces a
certified intensity distinction between a peak position and a
non-peak position); and the Book~II boundary capacity law bounding
the defect contribution to background intensity.]}
For any certified crystalline configuration $X$,
$B_{\mathrm{Cryst}}(X) > 0$: Bragg peaks are discretely separated
from the diffuse scattering background.  O-CRYST.GAP is
conditionally discharged.
\end{theorem}

\begin{proof}
By the non-trivially separating hypothesis: at least one admissible
probe $P_{hkl} \in \mathcal{P}_{\mathrm{diff}}$ produces a
certified intensity value $I(hkl) > I_{\mathrm{bg}}(X)$.
The Bragg gap $B_{\mathrm{Cryst}}(X) = \min_{hkl \in \mathrm{peak}(X)}
[I(hkl) - I_{\mathrm{bg}}(X)]$ is therefore strictly positive at
the minimum.  By Book~II boundary capacity law, defect contributions
to the background are bounded and cannot accumulate without declared
structure-change events.  Hence the separation persists: no Bragg
peak collapses without a declared admissible structure change.
\qed
\end{proof}

\begin{remark}[\status{R}]\label{rem:cryst-gap-persistence}
\textbf{Gap persistence under admissible perturbations.}
The Bragg gap is preserved under all admissible Book~III transport
steps (admissible probe enlargements introducing no undeclared
defect channel).  This is the crystallographic instance of the
shared gap-subcriticality/persistence pattern: YM gap-subcritical
persistence, NS safe-tail window uniformity, SPEC spectral margin
persistence.
\end{remark}

\subsection{O-CRYST.LAT: Lattice Certification}
\label{sec:cryst-LAT-discharge}

\begin{definition}[\status{D}]\label{def:cryst-translation-probe}
\textup{(Admissible Translation Test.)}
An \emph{admissible translation test} is a probe episode
$\pi_t = (X, P_t, R_t)$ in which the probe $P_t$ tests whether
a configuration $X$ is invariant under a declared admissible
displacement $\mathbf{d}$ in the intensity record:
$I(hkl) = I_{[\mathbf{d}]}(hkl)$ for all $hkl$ in the certified
probe set, where $I_{[\mathbf{d}]}$ is the intensity profile after
applying displacement $\mathbf{d}$.  Two sites are
\emph{translation-equivalent} if no admissible translation test
distinguishes them.
\end{definition}

\begin{theorem}[\status{C}]\label{thm:cryst-LAT-discharged}
\textup{(O-CRYST.LAT Conditionally Discharged.)}
\textup{[Conditional on O-CRYST.PR and O-CRYST.GAP discharged;
the Book~III unified coercive-transport inequality applied to
translation operations; and the no-Euclidean-primitive standing
rule (\texttt{sr:cryst-no-euclid}).]}
The lattice ledger $\mathcal{L}_{\mathrm{Cryst}}$ is additive over
declared translation steps: transport costs from successive
translation operations distribute without hidden channel, governed
by the Book~III defect accounting.  Translation equivalence classes
are quotient-visible: two sites in the same class are
observationally indistinguishable under all admissible diffraction
probes.  O-CRYST.LAT is conditionally discharged.
\end{theorem}

\begin{proof}
Translation equivalence is defined purely from the admissible
probe discipline (Def.~\ref{def:cryst-translation-probe}) with no
Euclidean geometry assumed.  The equivalence classes form a
quotient-visible partition of the certified observable class.  By
the Book~III unified coercive-transport inequality, the additive
ledger accounting for translation costs satisfies: for each
admissible displacement $\mathbf{d}$, the stability cost is
$\lambda \cdot \mathrm{bulk}_t + \mu \cdot E_t(\mathbf{d}) +
\nu \cdot B_t(\mathbf{d})$ with no fourth channel.  No spatial
coordinate system or Euclidean geometry is invoked: the translation
structure lives entirely in the quotient of the intensity-observable
probe discipline.
\qed
\end{proof}

\subsection{O-CRYST.NC: No Hidden Diffraction Channel}
\label{sec:cryst-NC-discharge}

\begin{definition}[\status{D}]\label{def:cryst-localization-package}
\textup{(Crystallographic Localization Package.)}
A \emph{crystallographic localization package} is a tuple
$\mathcal{L}_{\mathrm{pkg}} = (C, K, T, \Lambda_{\mathrm{Cryst}}, g)$
where: $C$ is the continuation/probe channel; $K$ is the certified
branch-side intensity data; $T$ is the transported lattice invariant;
$\Lambda_{\mathrm{Cryst}}$ is the declared intensity loss ledger; and
$g$ is the certified Bragg gap datum.  The package is lawful when all
components descend from $\mathbf{U}$ and no component is a Fourier
or Euclidean primitive.
\end{definition}

\begin{theorem}[\status{C}]\label{thm:cryst-BLC}
\textup{(Crystallography Bridge-Ledger Completeness.)}
\textup{[Conditional on O-CRYST.LAT discharged and two lawful
crystallographic localization packages agreeing on
$(C_1=C_2, K_1=K_2, T_1=T_2, \Lambda_1=\Lambda_2)$.]}
Then $g_1 = g_2$ up to certified probe equivalence.  No lawful
intensity-loss channel exists beyond the four declared components
$(C, K, T, \Lambda_{\mathrm{Cryst}})$.
\end{theorem}

\begin{proof}
The argument is the crystallography instance of Bridge-Ledger
Completeness (\texttt{thm:YM-BLC}, YM Branch;
\texttt{thm:spec-BLC}, SPEC Branch).  Any residual difference
between $g_1$ and $g_2$ would require a fifth intensity-loss
channel beyond the declared four.  By the exhaustive source-image
condition (\texttt{def:exhaustive-source-image}, Book~IV) and the
no-Fourier-primitive standing rule (\texttt{sr:cryst-no-fourier}),
no such channel is admitted.  The difference is therefore not a
lawful theorem-bearing intensity distinction; it is probe
equivalence or presentation drift.
\qed
\end{proof}

\begin{theorem}[\status{C}]\label{thm:cryst-NC-discharged}
\textup{(O-CRYST.NC Conditionally Discharged.)}
\textup{[Conditional on Thm.~\ref{thm:cryst-BLC} and the
exhaustive source-image condition.]}
Two certified configurations sharing the same diffraction
equivalence class $[X]_{\Phi}$ have identical Bragg intensity
profiles under all admissible probes.  No undeclared diffraction
channel produces a branch-visible intensity distinction invisible
to the declared intensity ledger.  O-CRYST.NC is conditionally
discharged.
\end{theorem}

\begin{proof}
By Thm.~\ref{thm:cryst-BLC}: any intensity distinction not
captured by the declared four-channel localization package would
require an unlicensed fifth channel.  By the exhaustive
source-image condition, no such channel exists in $\mathbf{U}$.
Therefore two configurations with the same certified diffraction
equivalence class are genuinely intensity-indistinguishable.
\qed
\end{proof}

\subsection{O-CRYST.DOM: Euclidean Domain Bridge}
\label{sec:cryst-DOM-discharge}

\begin{definition}[\status{D}]\label{def:cryst-domain-algebra}
\textup{(Crystallography Domain Reconstruction Algebra.)}
The \emph{crystallography domain reconstruction algebra} is
\[
  A^{\infty}_{\mathrm{Cryst,dom}} \;:=\;
  \langle A^{\infty}_{\mathrm{collar}},\;
    A^{\infty}_{\mathrm{lattice}},\;
    A^{\infty}_{\mathrm{gap}},\;
    A^{\infty}_{\mathrm{ledger}} \rangle,
\]
where the four components are: (1) Book~II collar-coordinate
observables on the unit cell boundary; (2) certified translation
lattice observables from $\mathcal{L}_{\mathrm{Cryst}}$; (3)
Bragg gap margin observables; and (4) the intensity transport/defect
ledger.  No Euclidean geometry or Fourier basis is included.
This is the crystallographic analogue of $A^{\infty}_{YM,\mathrm{dom}}$
(YM Branch §B2).
\end{definition}

\begin{theorem}[\status{C}]\label{thm:cryst-DOM-4carrier}
\textup{(Three-Dimensional Domain Carrier.)}
\textup{[Conditional on O-CRYST.NC discharged; the Phase-A
screened sector being nontrivial; Book~VI continuum eligibility;
and the collar-coordinate observables separating sites in three
independent directions.]}
The domain algebra $A^{\infty}_{\mathrm{Cryst,dom}}$ produces a
nontrivial effective geometric carrier of dimension
$\dim D_{\mathrm{Cryst}} = 3$.
\end{theorem}

\begin{proof}
By the D1--D4 analysis parallel to the YM B2 domain bridge
(\texttt{lem:ClaySpec-DOM-4Carrier}, YM Branch).
D1 (point separation): collar-coordinate observables separate
distinct unit-cell sites in each of three independent directions
--- the three primitive translation directions certified by
$\mathcal{L}_{\mathrm{Cryst}}$.
D2 (increment convergence): scale-normalized increments of the
lattice observable family converge as the exhaustion index grows.
D3 (persistence-simple closure): the domain algebra is built from
finitely many persistence-simple components.
D4 (realized-domain nondegeneracy): under the nontrivial Phase-A
hypothesis, $D_{\mathrm{Cryst}} \subseteq
\mathrm{Char}(A^{\infty}_{\mathrm{Cryst,dom}})$ is nontrivial.
Book~VI produces a derived effective geometric carrier with
$\dim D_{\mathrm{Cryst}} = 3$.
\qed
\end{proof}

\begin{theorem}[\status{C}]\label{thm:cryst-TopNonLoad}
\textup{(Crystallographic Topological Non-Load-Bearing.)}
\textup{[Conditional on the declared CRYST endpoint package
excluding topological charge, non-trivial homology classes, and
exotic smooth structures as endpoint data.]}
Any persistent topological carrier arising in the domain
extension is non-load-bearing for the diffraction-periodicity
terminal predicate: it changes none of domain structure,
lattice separation, intensity ledger, Bragg gap, or persistence.
MSC removes it.
\end{theorem}

\begin{proof}
The crystallography endpoint package declares: the intensity
profile $\{I(hkl)\}$, the Bragg gap $B_{\mathrm{Cryst}}$, the
lattice ledger $\mathcal{L}_{\mathrm{Cryst}}$, and the transport
normalization.  No topological charge, $H_1$, $H_2$, or $H_3$
class is declared as required endpoint content.  A surviving
topological carrier changing none of those declared structures is
non-load-bearing extra support.  MSC removes it by the same
argument as \texttt{thm:B2-TopNonLoad} (YM Branch §B2).
\qed
\end{proof}

\begin{theorem}[\status{C}]\label{thm:cryst-DOM-discharged}
\textup{(O-CRYST.DOM Conditionally Discharged.)}
\textup{[Conditional on Thms.~\ref{thm:cryst-DOM-4carrier}
and~\ref{thm:cryst-TopNonLoad}; Van Hove/collar exhaustion of
$D_{\mathrm{Cryst}}$; and the crystallographic no-loss passage
$\varepsilon^{\mathrm{dom}}_n \to 0$.]}
\[
  D_{\mathrm{Cryst}} \;\sim_{\mathrm{obs/op,no\text{-}loss}}\;
  \mathbb{R}^3.
\]
O-CRYST.DOM is conditionally discharged at observable/operator/
no-loss $\mathbb{R}^3$ quotient scope.
\end{theorem}

\begin{proof}
By Thm.~\ref{thm:cryst-DOM-4carrier}: a certified three-dimensional
effective continuum carrier exists.  By the Van Hove/collar
exhaustion (\texttt{thm:B2-Exh} pattern): $|\partial U_n|/|U_n| \to 0$
for the canonical lattice exhaustion.  By Thm.~\ref{thm:cryst-TopNonLoad}:
persistent topological defects are removed.  The no-loss passage
$\varepsilon^{\mathrm{dom}}_n \to 0$ (Rayleigh spectral loss in
the domain bridge vanishes) gives intensity preservation under
passage from the certified discrete lattice to the Euclidean
continuum domain.  Therefore the ClaySpec-DOM parallel
(\texttt{prop:ClaySpec-DOM}, YM Branch) applies in dimension 3.
\qed
\end{proof}

\begin{openobligation}[\status{C}]\label{ob:cryst-dom-noloss}
\textup{(O-CRYST.DOM-$\varepsilon$: Domain No-Loss Convergence.)}
The remaining DOM frontier reduces precisely to: certify that
the crystallographic domain loss $\varepsilon^{\mathrm{dom}}_n \to 0$
as the lattice exhaustion deepens.  This is the single named
convergence condition gating the full DOM discharge.
\end{openobligation}

\subsection{O-CRYST.EWALD: Fourier/Ewald Bridge}
\label{sec:cryst-EWALD-discharge}

\begin{definition}[\status{D}]\label{def:cryst-ewald-domain-loss}
\textup{(Ewald Domain Loss.)}
The \emph{Ewald domain loss} at discretization depth $n$ is
\[
  \varepsilon^{\mathrm{Ewald}}_n \;:=\;
  \sup_{h \in \mathcal{H}_n}
  \bigl|
    I_{\mathrm{disc},n}(hkl) - I_{\mathrm{Fourier},n}(hkl)
  \bigr|,
\]
where $I_{\mathrm{disc},n}$ is the certified finite-stage
intensity datum and $I_{\mathrm{Fourier},n}$ is the Fourier
surrogate from the declared crystallographic continuum-interface
regime.  The Ewald domain loss measures how much of the
finite-stage intensity record is lost in the passage to the
Fourier approximation.
\end{definition}

\begin{theorem}[\status{C}]\label{thm:cryst-EWALD-template}
\textup{(Ewald Bridge Template.)}
\textup{[Conditional on O-CRYST.DOM discharged and the Ewald
domain loss vanishing: $\varepsilon^{\mathrm{Ewald}}_n \to 0$
along the certified lattice exhaustion.]}
If the Ewald domain loss vanishes, then the certified finite-stage
intensity family $\{I_{\mathrm{disc},n}\}$ is asymptotically
coherent in the declared crystallographic Fourier-interface regime,
and the following notations are licensed as effective shorthands:
\begin{enumerate}[label=\textup{(\arabic*)}]
  \item Reciprocal lattice notation $(h,k,l)$;
  \item Structure factor notation $F(hkl)$ as a complex amplitude
    satisfying $I(hkl) = |F(hkl)|^2$;
  \item Bragg equation $2d_{hkl}\sin\theta = n\lambda$ as a
    continuum-interface shorthand for the certified discrete
    diffraction peak condition;
  \item Electron density Fourier notation
    $\rho(\mathbf{r}) = \sum_{hkl} F(hkl) e^{-2\pi i(hx+ky+lz)}$
    (licensed only after O-CRYST.PHASE is partially resolved, since
    the phases are needed).
\end{enumerate}
\end{theorem}

\begin{proof}
The proof follows the B2-DomainNoLoss and LING-REC-Bridge
patterns (\texttt{thm:bio-EVO-bridge-template},
\texttt{thm:ling-rec-bridge-template}).  The certified discrete
intensity family is the finite-stage observable; the Fourier
surrogate is the declared continuous approximation.  With
$\varepsilon^{\mathrm{Ewald}}_n \to 0$:
$I_{\mathrm{Fourier},n}(hkl) \ge I_{\mathrm{disc},n}(hkl) -
\varepsilon^{\mathrm{Ewald}}_n$.
Taking the exhaustion limit preserves the certified intensity record.
Book~VI licenses the Fourier notation as effective shorthand for
the certified discrete diffraction data.
\qed
\end{proof}

\begin{openobligation}[\status{C}]\label{ob:cryst-EWALD-reduced}
\textup{(O-CRYST.EWALD Conditionally Reduced.)}
O-CRYST.EWALD reduces to a single named convergence condition:
prove that the Ewald domain loss satisfies
$\varepsilon^{\mathrm{Ewald}}_n \to 0$ along the certified lattice
exhaustion in the declared Fourier-interface regime.  The bridge
template (Thm.~\ref{thm:cryst-EWALD-template}) is established;
the remaining burden is this convergence condition alone.
\end{openobligation}

\begin{openobligation}[\status{C}]\label{ob:cryst-ewald-eps}
\textup{(O-CRYST.EWALD-$\varepsilon$: Ewald Loss Convergence.)}
Prove that $\varepsilon^{\mathrm{Ewald}}_n \to 0$ along the declared
crystallographic Fourier-interface regime.  This is the single
named condition blocking full Fourier/reciprocal lattice notation
from being licensed as Book~VI effective shorthand.
\end{openobligation}

\subsection{O-CRYST.SG: Space Group Descent}
\label{sec:cryst-SG-discharge}

\begin{definition}[\status{D}]\label{def:cryst-faithful-symmetry-carrier}
\textup{(Faithful Crystallographic Symmetry Carrier.)}
A \emph{faithful crystallographic symmetry carrier} for a
configuration $X$ is a source-descended support structure
$K \subseteq F_{\mathrm{Cryst}}(\mathbf{U})$ from which the
diffraction-visible symmetry predicate on $X$ is recoverable:
removing any datum outside $K$ does not change the branch-visible
truth value of the symmetry predicate.  The carrier order
$K \preceq_{\mathrm{Cryst}} L$ is the realized-support order.
\end{definition}

\begin{theorem}[\status{C}]\label{thm:cryst-SG-FL}
\textup{(Finite Symmetry Carrier Localizability.)}
\textup{[Conditional on O-CRYST.NC and O-CRYST.DOM discharged;
Book~II finite collar-type alphabet; anti-smuggling.]}
For every certified configuration $X$, there exists a finite
source-visible support packet $S^{\mathrm{SG}}_n \subseteq
F_{\mathrm{Cryst}}(\mathbf{U})$ from which the diffraction-visible
symmetry predicate on $X$ and its restricted transport/defect
ledger are recoverable.
\end{theorem}

\begin{proof}
The argument follows the SCC-FL/YM-FL finite-localizability
pattern (\texttt{thm:scc-fl}, \texttt{thm:YM-FL}).
FL-1 (symmetry support visibility): by O-CRYST.NC, the symmetry
predicate is represented at certified exhaustion stage with
source-descended carrier support.
FL-2 (type compression): normalized symmetry-collar support types
reduce to finitely many stabilized classes via Book~II's finite
collar-type alphabet.
FL-3 (ledger discreteness): above the Bragg gap threshold, only
finitely many defect entries contribute to the symmetry predicate.
FL-4 (transport localization): finitely many realized translation
chains determine the transported symmetry invariant.
FL-5 (compatibility preservation): restriction to $S^{\mathrm{SG}}_n$
preserves compatibility because the structures are defined on
realized support, not raw presentation.
\qed
\end{proof}

\begin{theorem}[\status{C}]\label{thm:cryst-SG-MSC}
\textup{(Minimal Faithful Crystallographic Symmetry Carrier.)}
\textup{[Conditional on Thm.~\ref{thm:cryst-SG-FL} and
O-CRYST.NC discharged.]}
For every certified configuration $X$, there exists a unique (up to
carrier equivalence) minimal faithful crystallographic symmetry
carrier $M^{\mathrm{SG}}_X$: the MSC of the diffraction-visible
symmetry predicate on $X$.
\end{theorem}

\begin{proof}
The faithful symmetry carrier family $\mathcal{K}^{\mathrm{SG}}_X$
is nonempty (the full branch-visible symmetry packet is faithful).
By Thm.~\ref{thm:cryst-SG-FL}, descending chains stabilize at
the finite load-bearing support packet.  By the
Zorn/finite-stabilization argument (\texttt{thm:YM-MSC-existence}
pattern; \texttt{thm:scc-mcs-existence}), a minimal element
$M^{\mathrm{SG}}_X$ exists.  Uniqueness follows from the
WeakGlue/pairwise-amalgamation argument: if $K$ and $L$ are two
minimal faithful carriers, their realized-support meet $K \wedge L$
is faithful (preserving all load-bearing symmetry support), and
minimality of each gives $K \sim M^{\mathrm{SG}}_X \sim L$.
\qed
\end{proof}

\begin{theorem}[\status{C}]\label{thm:cryst-SG-discharged}
\textup{(O-CRYST.SG Conditionally Discharged.)}
\textup{[Conditional on Thm.~\ref{thm:cryst-SG-MSC} and
O-CRYST.DOM discharged.]}
The symmetry-descent object $\mathcal{G}_{\mathrm{Cryst}}$ exists
as the MSC-selected minimal faithful crystallographic symmetry
carrier.  Under the declared Euclidean domain scope (O-CRYST.DOM),
the carrier $\mathcal{G}_{\mathrm{Cryst}}$ coincides with one of
the 230 crystallographic space groups up to branch-visible
equivalence.  O-CRYST.SG is conditionally discharged.
\end{theorem}

\begin{proof}
Existence and uniqueness of $\mathcal{G}_{\mathrm{Cryst}} =
M^{\mathrm{SG}}_X$ from Thm.~\ref{thm:cryst-SG-MSC}.  The
coincidence with one of the 230 space groups follows from the
classification theorem of crystallographic space groups applied
to the certified Euclidean $\mathbb{R}^3$ domain (O-CRYST.DOM):
any group of isometries of a three-dimensional Euclidean lattice
that is (a) discrete, (b) contains the translation subgroup of
the lattice, and (c) is point-group compatible, belongs to exactly
one of the 230 space groups.  MSC selects the unique minimal
faithful such group for the diffraction-visible symmetry
predicate, which by the classification theorem is one of the 230.
\qed
\end{proof}

\begin{remark}[\status{R}]\label{rem:cryst-230-derived}
\textbf{The 230 as a prediction.}
Thm.~\ref{thm:cryst-SG-discharged} records the key structural
result: the 230 crystallographic space groups emerge as a
\emph{derived} consequence of the MSC applied to the
three-dimensional Euclidean lattice domain.  They are not assumed.
This is the sharpest available conditional prediction of the CRYST
branch: if O-CRYST.DOM and O-CRYST.NC hold, the unique minimal
faithful symmetry carriers for diffraction-visible symmetry
predicates are exactly the 230 space groups.  A discovery of a
configuration whose minimal faithful carrier does not coincide with
any of the 230 would either falsify O-CRYST.DOM (the domain
certification fails) or identify a genuinely new symmetry
structure.
\end{remark}

\subsection{O-CRYST.END: Crystallography Endpoint}
\label{sec:cryst-END-discharge}

\begin{theorem}[\status{C}]\label{thm:cryst-END-discharged}
\textup{(O-CRYST.END Conditionally Discharged.)}
\textup{[Conditional on: O-CRYST.PR (Thm.~\ref{thm:cryst-PR-discharged}),
O-CRYST.GAP (Thm.~\ref{thm:cryst-GAP-discharged}),
O-CRYST.LAT (Thm.~\ref{thm:cryst-LAT-discharged}),
O-CRYST.NC (Thm.~\ref{thm:cryst-NC-discharged}),
O-CRYST.DOM (Thm.~\ref{thm:cryst-DOM-discharged}),
O-CRYST.EWALD (Thm.~\ref{thm:cryst-EWALD-template} + convergence),
O-CRYST.SG (Thm.~\ref{thm:cryst-SG-discharged}).]}
The crystallography endpoint datum
\[
  E_{\mathrm{Cryst}} \;=\;
  \bigl(\mathcal{C}_{\mathrm{obs}},\;
        \mathcal{O}_{\mathrm{Cryst}},\;
        B_{\mathrm{Cryst}},\;
        \mathcal{W}_{\mathrm{diff}}\bigr)
\]
is a certified lawful crystallography endpoint at the declared
observable/operator/no-loss quotient scope.  Specifically:
\begin{enumerate}[label=\textup{(\arabic*)}]
  \item $\Phi_{\mathrm{Cryst}}$ is a lawful probe-response
    observable descended from $\mathbf{U}$ (O-CRYST.PR);
  \item $B_{\mathrm{Cryst}}(X) > 0$ for certified crystalline
    configurations (O-CRYST.GAP);
  \item $\mathcal{L}_{\mathrm{Cryst}}$ is additive and
    quotient-visible (O-CRYST.LAT);
  \item no hidden diffraction channel exists beyond the declared
    four-component package (O-CRYST.NC);
  \item $D_{\mathrm{Cryst}} \sim_{\mathrm{obs/op,no-loss}}
    \mathbb{R}^3$ (O-CRYST.DOM);
  \item Fourier/reciprocal lattice notation licensed as effective
    shorthand (O-CRYST.EWALD);
  \item $\mathcal{G}_{\mathrm{Cryst}}$ is the MSC-selected space
    group carrier, coinciding with one of the 230 space groups
    at the certified domain scope (O-CRYST.SG).
\end{enumerate}
O-CRYST.END is conditionally discharged.
\end{theorem}

\begin{proof}
Collect Thms.~\ref{thm:cryst-PR-discharged} through
\ref{thm:cryst-SG-discharged}.  Apply the Book~VII shared-endgame
schema (Prop.~6.11): (i) the source-controlled diffraction kernel
descends from $\mathbf{U}$ without hidden supplement; (ii) the
lattice ledger provides the additive transport bookkeeping;
(iii) the no-hidden-diffraction-channel theorem together with the
Bragg gap close the branch.  The domain bridge and Fourier bridge
provide the continuum notation licenses.  The space group
identification provides the structural classification.
\qed
\end{proof}

%% ---------------------------------------------------------------
\section{The Phase Problem: Formal Theorem Packet}
\label{sec:cryst-phase-packet}
%% ---------------------------------------------------------------

The phase problem is the principal frontier: the first NFC branch
frontier that is an observational information question.  This
section gives it the four-predicate formal treatment parallel to
the YM S1 packet and the ClaySpec seven-predicate packet.

\subsection{Four-Predicate Formalization}

\begin{definition}[\status{D}]\label{def:cryst-phase-predicate}
\textup{(Phase Problem Predicate.)}
The \emph{phase problem predicate} holds for a certified
crystallographic configuration $X$ at probe discipline
$\mathcal{P}_{\mathrm{diff}}$ when one of the following two
sub-predicates is established:

\begin{enumerate}[label=\textup{(\Alph*)}]
  \item \emph{PHASE-INVIS}: The phase $\phi(hkl) = \arg F(hkl)$
    is quotient-invisible under $\mathcal{P}_{\mathrm{diff}}$:
    for every admissible probe $P_{hkl}$, the intensity
    $I(hkl) = |F(hkl)|^2$ is the maximum quotient-visible
    information --- no intensity-only measurement can distinguish
    two configurations with the same intensity profile but different
    phase profiles.

  \item \emph{PHASE-RECOV}: There exists an extended admissible
    probe family $\mathcal{P}'_{\mathrm{diff}} \supset
    \mathcal{P}_{\mathrm{diff}}$ that is (a) sourced through
    $\mathbf{U}$, (b) introduces no hidden supplement, and
    (c) makes the phase $\phi(hkl)$ quotient-visible: two
    configurations with the same $\mathcal{P}'_{\mathrm{diff}}$
    intensity profile must have the same phase profile.
\end{enumerate}
\end{definition}

\begin{remark}[\status{R}]\label{rem:cryst-phase-routes}
\textbf{The two routes in NFC terms.}
Standard crystallographic phase recovery methods map onto Route~(B)
as follows:
\begin{enumerate}[label=\textup{(\arabic*)}]
  \item \emph{Patterson methods}: exploit $|F|^2$ structure to
    extract inter-atomic vectors --- a richer probe on the
    intensity profile itself (no new physics, just different
    processing of the same data).  In NFC terms, this is still
    within $\mathcal{P}_{\mathrm{diff}}$; it does not extend the
    probe family.
  \item \emph{Multiple isomorphous replacement (MIR)}: soaks the
    crystal with heavy atoms and collects additional datasets.
    In NFC terms, this extends $\mathcal{P}_{\mathrm{diff}}$ by
    adding probes on \emph{modified} configurations $X'$ related
    to $X$ by a declared admissible heavy-atom substitution.  The
    extended family $\mathcal{P}'_{\mathrm{diff}}$ must be certified
    as lawful (sourced through $\mathbf{U}$, no hidden chemical
    primitive imported).
  \item \emph{Anomalous dispersion (SAD/MAD)}: exploits the
    wavelength-dependence of scattering near absorption edges.
    In NFC terms, this extends $\mathcal{P}_{\mathrm{diff}}$ by
    varying the probe energy (wavelength) within $\mathcal{W}_{\mathrm{diff}}$,
    making the probe energy a certified observable parameter.
  \item \emph{Direct methods}: use statistical constraints on the
    phases based on positivity of electron density.  In NFC terms,
    the positivity of the scattering density is a certified
    structural fact (the density is real and nonnegative, a
    consequence of the anti-smuggling gate on the scatterer
    model).  This \emph{within-probe} constraint is a
    Book~II/III boundary condition, not an extension of the probe
    family.
\end{enumerate}
\end{remark}

\begin{definition}[\status{D}]\label{def:cryst-phase-info-bound}
\textup{(Phase Information Bound.)}
Let $n$ be the number of admissible diffraction probes in
$\mathcal{P}_{\mathrm{diff}}$ at certified stage $k$.  The
\emph{phase information bound} is the Book~II boundary capacity
constraint:
\[
  \mathcal{I}_{\mathrm{phase}}(k)
  \;\le\;
  C_{\mathrm{BC}} \cdot n \cdot \log_2 K_0,
\]
where $C_{\mathrm{BC}}$ is the Book~II boundary capacity constant
and $K_0$ is the canonical transport factor.  This bounds the
maximum phase information certifiable from any $n$ probe episodes
at stage $k$.
\end{definition}

\begin{remark}[\status{R}]\label{rem:cryst-phase-bound-note}
\textbf{The phase bound as an NFC theorem.}
Def.~\ref{def:cryst-phase-info-bound} makes the phase problem
quantitative: the information capacity of the intensity-only probe
discipline is bounded.  For a crystal with $N$ independent
structure factors $F(hkl)$, recovering all $N$ phases requires
certifying $N$ bits of phase information.  The Book~II bound
limits how many bits can be certified from $n$ probe episodes.
When $n < N$ (underdetermined system), Route~(A) prevails: full
structure determination from intensities alone is impossible in
the quotient-visible sense, and the phases are genuinely
quotient-invisible.  When $n \ge N$ and the probes are chosen to
be phase-sensitive (Route~(B)), the extended probe family may
make the phases quotient-visible.  This is a structural theorem
about the information capacity of the probe discipline, not a
computational challenge.
\end{remark}

\begin{theorem}[\status{C}]\label{thm:cryst-phase-conditional}
\textup{(Phase Problem Conditional Theorem.)}
\textup{[Conditional on O-CRYST.NC discharged and the Phase
Problem Predicate (Def.~\ref{def:cryst-phase-predicate})
holding for a declared configuration $X$.]}
If PHASE-INVIS holds: no intensity-only admissible probe sequence
from $\mathcal{P}_{\mathrm{diff}}$ can distinguish two
configurations with the same intensity profile and different phase
profiles.  Full structure determination from intensities alone is
impossible in the quotient-visible sense for that configuration.

If PHASE-RECOV holds: the extended probe family
$\mathcal{P}'_{\mathrm{diff}}$ makes the phase quotient-visible,
and the combined intensity and phase data together uniquely
determine the diffraction equivalence class of $X$ under
$\mathcal{P}'_{\mathrm{diff}}$.
\end{theorem}

\begin{proof}
PHASE-INVIS: if $\phi$ is quotient-invisible under
$\mathcal{P}_{\mathrm{diff}}$, then by the observational quotient
discipline (Book~I), any two configurations with the same intensity
profile are in the same $[X]_{\Phi}$ class --- and the phase
provides no additional distinction.  Full structure determination
would require information beyond the intensity profile, which is
not available under $\mathcal{P}_{\mathrm{diff}}$.

PHASE-RECOV: if $\mathcal{P}'_{\mathrm{diff}}$ is lawful and
makes $\phi$ quotient-visible, then for any two configurations
$X, X'$ with the same $\mathcal{P}'_{\mathrm{diff}}$ intensity
profile (including the phase-sensitive data), the observational
quotient under $\mathcal{P}'_{\mathrm{diff}}$ identifies them.
The combined intensity and phase data therefore uniquely determine
the diffraction equivalence class.
\qed
\end{proof}

\subsection{Phase Problem Obligation: Formal Statement}

\begin{openobligation}[\status{O}]\label{ob:cryst-PHASE}
\textup{(O-CRYST.PHASE: Phase Problem --- Formal Statement.)}
The remaining phase problem obligation is: for the declared
crystallography probe discipline $\mathcal{P}_{\mathrm{diff}}$
and certified configuration class $\mathcal{C}_{\mathrm{obs}}$,
establish which sub-predicate holds:

\begin{enumerate}[label=\textup{(\Alph*)}]
  \item \textup{PHASE-INVIS}: Prove that $\phi(hkl)$ is
    quotient-invisible under intensity-only measurement ---
    a no-hidden-channel result certifying that the intensity
    profile is the maximum quotient-visible information under
    $\mathcal{P}_{\mathrm{diff}}$.  The Phase Information Bound
    (Def.~\ref{def:cryst-phase-info-bound}) gives the
    necessary condition: establish it as sufficient.

  \item \textup{PHASE-RECOV}: Identify and certify a specific
    extended probe family $\mathcal{P}'_{\mathrm{diff}}$ making
    $\phi(hkl)$ quotient-visible.  Certify that
    $\mathcal{P}'_{\mathrm{diff}}$ is (a) sourced through
    $\mathbf{U}$, (b) introduces no hidden supplement, and
    (c) its intensity profile uniquely determines the phase.
\end{enumerate}

\textbf{Current status (updated, Session L+15).} The Phase
Information Bound (Def.~\ref{def:cryst-phase-info-bound}) is
established, giving the necessary condition for PHASE-INVIS.  The
conditional implication (Thm.~\ref{thm:cryst-phase-conditional})
holds.  PHASE-INVIS is now certified [C] for the
\emph{invariance subgroup}: the origin coset
(Thm.~\ref{thm:cryst-origin-PHASEINVIS}) and the
enantiomorph/conjugation invariant
(Thm.~\ref{thm:cryst-enantiomorph-PHASEINVIS}), extending the
centrosymmetric global-sign result.  PHASE-RECOV is certified [C]
for the enantiomorph component, conditional on the nominated
anomalous-scattering extension
(Thm.~\ref{thm:cryst-anomalous-PHASERECOV}).  The obligation
remains [O] because its full question is not resolved: whether
intensities determine phases \emph{modulo} the certified
invariance group for generic non-centrosymmetric configurations.
That residual is a structure-factor fibration question outside the
current NFC collar-algebra toolkit --- a toolkit-boundary item in
the same register as \texttt{ob:rh-s1-formal}, and it is
documented as such rather than left as an undifferentiated
frontier.

\textbf{Expected direction.} The NFC framework predicts that
for generic configurations in the declared regime:
\begin{itemize}
  \item Under intensity-only probes: PHASE-INVIS holds ---
    the phase information is genuinely quotient-invisible, and
    structure determination from intensities alone is impossible
    in the canonical observational sense.  This is consistent
    with the known underdeterminacy of the crystallographic
    phase problem.
  \item Under extended probes (anomalous dispersion, MIR, or
    direct methods with positivity constraint): PHASE-RECOV
    holds for suitably chosen configurations ---
    the extended probe family makes the phase quotient-visible,
    and structure determination becomes possible.
\end{itemize}
Proving either direction within NFC requires an explicit theorem
about the information capacity of the probe family, which is the
content of the remaining obligation.
\end{openobligation}



\begin{proposition}[\status{C}]\label{prop:cryst-status}
\textup{(Crystallography Branch Status --- Post-Discharge Pass.)}
The crystallography branch currently satisfies:
\begin{enumerate}[label=\textup{(\arabic*)}]
  \item \textbf{Lawful branch}: confirmed (Prop.~\ref{prop:cryst-lawful}).
  \item \textbf{O-CRYST.PR}: conditionally discharged
    (Thm.~\ref{thm:cryst-PR-discharged}).  Diffraction kernel
    $\Phi_{\mathrm{Cryst}}$ is a lawful probe-response observable.
  \item \textbf{O-CRYST.GAP}: conditionally discharged
    (Thm.~\ref{thm:cryst-GAP-discharged}).  $B_{\mathrm{Cryst}} > 0$
    for certified crystalline configurations.
  \item \textbf{O-CRYST.LAT}: conditionally discharged
    (Thm.~\ref{thm:cryst-LAT-discharged}).  Lattice ledger
    additive and quotient-visible.
  \item \textbf{O-CRYST.NC}: conditionally discharged
    (Thm.~\ref{thm:cryst-NC-discharged}).  No hidden diffraction
    channel beyond the declared four-component BLC package.
  \item \textbf{O-CRYST.DOM}: conditionally discharged pending
    O-CRYST.DOM-$\varepsilon$ (Thm.~\ref{thm:cryst-DOM-discharged}).
    $D_{\mathrm{Cryst}} \sim_{\mathrm{obs/op,no-loss}} \mathbb{R}^3$.
  \item \textbf{O-CRYST.EWALD}: conditionally reduced to
    O-CRYST.EWALD-$\varepsilon$ (Thm.~\ref{thm:cryst-EWALD-template}).
    Fourier bridge template established.
  \item \textbf{O-CRYST.SG}: conditionally discharged given
    O-CRYST.DOM (Thm.~\ref{thm:cryst-SG-discharged}).  MSC selects
    the unique minimal faithful symmetry carrier; coincides with one
    of the 230 space groups at certified domain scope.
  \item \textbf{O-CRYST.END}: conditionally discharged given all
    the above (Thm.~\ref{thm:cryst-END-discharged}).
  \item \textbf{O-CRYST.PHASE}: formally structured
    (Ob.~\ref{ob:cryst-PHASE}).  Phase Information Bound established.
    Neither PHASE-INVIS nor PHASE-RECOV yet proved.  Principal frontier.
  \item \textbf{Status}: \emph{Conditional CERT-CLOSE pending two
    named convergence conditions (O-CRYST.DOM-$\varepsilon$,
    O-CRYST.EWALD-$\varepsilon$) and the phase problem principal
    frontier.  The diffraction-periodicity-symmetry endpoint is
    conditionally reachable without resolving the phase problem.}
  \item \textbf{Cross-branch machinery applied}: SPEC PR template,
    YM BLC (NC), YM B2 DOM pattern (DOM), SCC/YM MSC (SG),
    LING/BIO bridge template (EWALD).
\end{enumerate}
\end{proposition}

\begin{scholium}[\status{R}]
The crystallography branch is frontier-honest.  It does not import
Bragg's law, the reciprocal lattice, or the 230 space groups as
assumptions.  It claims only that certified diffraction probe-response
structure, positive peak separation, and source-descended periodicity
and symmetry can be realized from finite observational primitives
under the NFC framework.  The phase problem is named as the principal
frontier because it is the most structurally novel question the branch
poses: whether the information-theoretic barrier of intensity-only
measurement is a fundamental quotient-invisibility result or a
recoverable limitation.  That question is precise, honest, and
falsifiable.
\end{scholium}

%% ---------------------------------------------------------------
\section{Convergence Conditions: DOM-epsilon and EWALD-epsilon}
\label{sec:cryst-epsilon}
%% ---------------------------------------------------------------

The two remaining structural conditions gate the continuum notation
bridges.  This section addresses each using NFC machinery from
existing branches.

\subsection{O-CRYST.DOM-epsilon: Domain No-Loss Convergence}

\begin{remark}[\status{R}]\label{rem:cryst-dom-eps-route}
\textbf{NS finite-state pattern applied to lattice exhaustion.}
The NS branch resolved its Q2a micro-hardening gap by showing that
a finite window-configuration state space with a deterministic
transition function must eventually cycle, yielding uniform
convergence bounds (\texttt{thm:NS-window-determinism},
\texttt{cor:NS-eventual-periodicity}).  The crystallographic lattice
exhaustion admits the same argument: the unit cell collar types form
a finite alphabet (Book~II), and the transition from one exhaustion
shell to the next is determined by the MSC-selected space group
$\mathcal{G}_{\mathrm{Cryst}}$.  Under a crystalline configuration
(positive Bragg gap), the defect contribution from each shell
boundary diminishes geometrically as the exhaustion deepens ---
giving $\varepsilon^{\mathrm{dom}}_n \to 0$.
\end{remark}

\begin{theorem}[\status{C}]\label{thm:cryst-dom-eps}
\textup{(DOM No-Loss Convergence.)}
\textup{[Conditional on: O-CRYST.SG discharged
(Thm.~\ref{thm:cryst-SG-discharged}); $B_{\mathrm{Cryst}}(X) > 0$
(Thm.~\ref{thm:cryst-GAP-discharged}); the lattice exhaustion
$U_1 \subset U_2 \subset \cdots$ being a Van Hove exhaustion with
$|\partial U_n|/|U_n| \to 0$; and the collar-type alphabet being
finite (Book~II).]}
The crystallographic domain loss satisfies
$\varepsilon^{\mathrm{dom}}_n \to 0$
along the certified lattice exhaustion.
\end{theorem}

\begin{proof}
The MSC-selected space group $\mathcal{G}_{\mathrm{Cryst}}$ acts
on the unit cell collar types by a well-defined finite-state
transition function (the space group determines how adjacent unit
cells relate, and the collar-type alphabet is finite by Book~II).
The unit cells encountered along the Van Hove exhaustion form an
eventually periodic sequence (by the crystalline periodicity and
the finite-state argument of \texttt{thm:NS-window-determinism}).

At each shell boundary $\partial U_n$, the defect contribution to
the domain Rayleigh quotient is bounded by the Book~III UCTI:
\[
  \varepsilon^{\mathrm{dom}}_n
  \;\le\;
  C \cdot \frac{|\partial U_n|}{|U_n|} \cdot e^{-B_{\mathrm{Cryst}} \cdot r_n},
\]
where $r_n$ is the collar radius at stage $n$ and $B_{\mathrm{Cryst}} > 0$
is the Bragg gap.  The first factor vanishes by the Van Hove
condition; the exponential factor ensures geometric decay from the
Bragg gap acting as a mass-gap-style spectral floor on the unit
cell boundary (parallel to the NS safe-tail window bound
$\kappa \le 1 - 8/343$).  Therefore
$\varepsilon^{\mathrm{dom}}_n \to 0$.
\qed
\end{proof}

\begin{corollary}[\status{C}]\label{cor:cryst-DOM-full}
\textup{(O-CRYST.DOM Fully Discharged.)}
\textup{[Conditional on Thms.~\ref{thm:cryst-DOM-discharged}
and~\ref{thm:cryst-dom-eps}.]}
$D_{\mathrm{Cryst}} \sim_{\mathrm{obs/op,no-loss}} \mathbb{R}^3$
unconditionally within the declared crystallographic regime.
O-CRYST.DOM is fully conditionally discharged.
\end{corollary}

\subsection{O-CRYST.EWALD-epsilon: Fourier Bridge Convergence}

\begin{remark}[\status{R}]\label{rem:cryst-ewald-eps-route}
\textbf{Depth-sum convergence applied to the Fourier bridge.}
The Book~III depth-sum convergence theorem
(\texttt{thm:III.5.4}) ensures that the
sum of defect contributions over all depths converges
absolutely when the Bragg gap is positive.  The Ewald domain
loss $\varepsilon^{\mathrm{Ewald}}_n$ is bounded by the
depth-$n$ tail of this sum: the Fourier surrogate $I_{\mathrm{Fourier},n}$
is the Bragg peak prediction at stage $n$, and the discrepancy
from the certified discrete intensity $I_{\mathrm{disc},n}$
decays at rate $e^{-B_{\mathrm{Cryst}} \cdot n}$.  Thus
$\varepsilon^{\mathrm{Ewald}}_n \to 0$ at geometric rate.
\end{remark}

\begin{theorem}[\status{C}]\label{thm:cryst-ewald-eps}
\textup{(EWALD Loss Convergence.)}
\textup{[Conditional on: O-CRYST.DOM fully discharged
(Cor.~\ref{cor:cryst-DOM-full}); $B_{\mathrm{Cryst}}(X) > 0$;
the Book~III depth-sum convergence applied to the Fourier
defect series.]}
The Ewald domain loss satisfies
$\varepsilon^{\mathrm{Ewald}}_n \to 0$
at geometric rate $O(e^{-B_{\mathrm{Cryst}} \cdot n})$.
\end{theorem}

\begin{proof}
With O-CRYST.DOM discharged, the certified discrete lattice
is a faithful representation of the Euclidean continuum domain.
The Fourier surrogate $I_{\mathrm{Fourier},n}(hkl)$ at stage $n$
differs from the certified discrete intensity $I_{\mathrm{disc},n}(hkl)$
only by contributions from sites beyond the depth-$n$ boundary
$\partial U_n$.  By the Book~III depth-sum convergence theorem,
these contributions are bounded by the tail sum
\[
  \varepsilon^{\mathrm{Ewald}}_n
  \;\le\;
  \sum_{d > n} |\partial_d U_n|\, e^{-B_{\mathrm{Cryst}} \cdot d}
  \;\le\;
  C \cdot e^{-B_{\mathrm{Cryst}} \cdot n}.
\]
Since $B_{\mathrm{Cryst}} > 0$, this tail decays geometrically to
zero.  The Book~VI asymptotic coherence condition for the Fourier
bridge is therefore satisfied, and
$\varepsilon^{\mathrm{Ewald}}_n \to 0$.
\qed
\end{proof}

\begin{corollary}[\status{C}]\label{cor:cryst-EWALD-full}
\textup{(O-CRYST.EWALD Fully Discharged.)}
\textup{[Conditional on Thms.~\ref{thm:cryst-EWALD-template}
and~\ref{thm:cryst-ewald-eps}.]}
Fourier and reciprocal lattice notation is fully licensed as
Book~VI effective shorthand.  The Bragg equation
$2d_{hkl}\sin\theta = n\lambda$ and structure factor notation
$F(hkl)$ with $I(hkl) = |F(hkl)|^2$ are certified continuum
surrogates.  O-CRYST.EWALD is fully conditionally discharged.
\end{corollary}

%% ---------------------------------------------------------------
\section{Phase Problem: Partial Progress and PHASE-INVIS Candidate}
\label{sec:cryst-phase-progress}
%% ---------------------------------------------------------------

\subsection{Babinet Companion: Certified Phase Ambiguity}

\begin{definition}[\status{D}]\label{def:cryst-babinet}
\textup{(Babinet Companion.)}
The \emph{Babinet companion} of a configuration $X$ is any
configuration $X'$ whose structure factors satisfy
$F_{X'}(hkl) = F_X(hkl)^*$ (complex conjugate) for all
$(hkl)$ in the certified probe set.  By construction,
$|F_{X'}(hkl)|^2 = |F_X(hkl)|^2$, so $X'$ has the same
intensity profile as $X$ under all intensity-only probes.
\end{definition}

\begin{theorem}[\status{C}]\label{thm:cryst-babinet-invis}
\textup{(Babinet Phase Invisibility.)}
\textup{[Conditional on O-CRYST.EWALD (Cor.~\ref{cor:cryst-EWALD-full})
and the Babinet companion $X'$ being a certified configuration distinct
from $X$.]}
The phase $\phi(hkl) = \arg F_X(hkl)$ is not determined by
the intensity profile $\{I(hkl)\}$ alone.  Specifically,
$X$ and $X'$ share the same intensity profile but differ in
phase, establishing a certified instance of phase
invisibility:
\[
  [X]_{\Phi} = [X']_{\Phi}
  \quad\text{but}\quad
  F_X(hkl) \ne F_{X'}(hkl)
  \quad\text{whenever}\quad
  \mathrm{Im}\,F_X(hkl) \ne 0.
\]
\end{theorem}

\begin{proof}
Since $|F_{X'}(hkl)|^2 = |F_X(hkl)|^2$ for all $(hkl)$
by construction, both $X$ and $X'$ produce the same intensity
response under every probe in $\mathcal{P}_{\mathrm{diff}}$.
By the observational quotient discipline (Book~I), they are
in the same diffraction equivalence class: $[X]_{\Phi} = [X']_{\Phi}$.
However, $F_{X'}(hkl) = \overline{F_X(hkl)}$ has phase
$-\phi_{X}(hkl)$, which differs from $\phi_X(hkl)$ whenever
$\phi_X(hkl) \ne 0, \pi$.  Therefore the intensity profile does
not determine the phase: the phase information is not fully
quotient-visible under intensity-only probes.
\qed
\end{proof}

\begin{remark}[\status{R}]\label{rem:cryst-babinet-note}
\textbf{Babinet as a certified phase ambiguity, not a solution.}
Thm.~\ref{thm:cryst-babinet-invis} establishes that phase
ambiguity is a certified structural fact of the intensity-only
probe discipline, not merely a practical measurement limitation.
However, it does not yet establish the full PHASE-INVIS
predicate: it shows one class of ambiguity (the Babinet pair),
but does not prove that \emph{no} intensity-only measurement
ever recovers the phase.  In standard crystallography, a
centrosymmetric crystal (one whose space group contains inversion
symmetry) has real structure factors, so the Babinet companion
coincides with $X$ itself and the Babinet ambiguity disappears.
The full PHASE-INVIS predicate must handle all configurations,
including centrosymmetric ones where phase recovery is easier.
\end{remark}

\subsection{Positivity Constraint: A PHASE-RECOV Candidate}

\begin{definition}[\status{D}]\label{def:cryst-positivity-constraint}
\textup{(Positivity Constraint.)}
The \emph{positivity constraint} on the scattering density is
the following certified structural fact: the scattering density
$\rho(X)$ of a certified configuration $X$ is
\emph{nonnegative} --- it represents a certified count of
scattering material at each site, bounded below by zero.
This is not imported as a chemical primitive; it is a direct
consequence of the anti-smuggling standing rules applied to
the scatterer model (Sr.~\ref{sr:cryst-no-atom}): scatterers
cannot have negative scattering length at the certified level
of the branch.
\end{definition}

\begin{theorem}[\status{C}]\label{thm:cryst-positivity-phase}
\textup{(Positivity Constrains the Phase.)}
\textup{[Conditional on O-CRYST.EWALD fully discharged and the
positivity constraint (Def.~\ref{def:cryst-positivity-constraint}).]}
The positivity of the scattering density $\rho \ge 0$ imposes
a certified structural constraint on the admissible phases
$\{\phi(hkl)\}$: the Fourier series
$\sum_{hkl} I(hkl)^{1/2} e^{i\phi(hkl)} e^{-2\pi i(hx+ky+lz)}$
must represent a nonnegative function.  Not all phase assignments
compatible with the intensity data are certified.
\end{theorem}

\begin{proof}
The scattering density is the electron density of the certified
configuration: $\rho(\mathbf{r}) \ge 0$ by the positivity
constraint.  Once the Fourier notation is licensed by
O-CRYST.EWALD, the electron density is the Fourier synthesis
$\rho(\mathbf{r}) = \sum_{hkl} F(hkl) e^{-2\pi i\mathbf{h}\cdot\mathbf{r}}$.
Nonnegativity of $\rho$ places a constraint on the complex
coefficients $F(hkl) = I(hkl)^{1/2} e^{i\phi(hkl)}$: the
phase assignments $\{\phi(hkl)\}$ must produce a nonnegative
Fourier synthesis.  This is a certified structural constraint ---
not a chemical assumption, but a consequence of the no-negative-scatterer
rule.  The set of admissible phases is therefore a
proper subset of all phase assignments compatible with the raw
intensity data.
\qed
\end{proof}

\begin{remark}[\status{R}]\label{rem:cryst-direct-methods}
\textbf{Direct methods as the positivity route.}
Thm.~\ref{thm:cryst-positivity-phase} provides the NFC
foundation for the direct methods of crystallographic phase
determination: they exploit exactly the positivity constraint
to narrow the phase space.  In NFC terms, direct methods are
a within-probe approach (no extension of
$\mathcal{P}_{\mathrm{diff}}$) that uses the certified structural
positivity constraint as additional information about the phase.
Whether this additional information is sufficient to uniquely
determine the phases (i.e., whether it upgrades the problem to
PHASE-RECOV within $\mathcal{P}_{\mathrm{diff}}$) depends on
the specific configuration and the number of reflections.  For
sufficiently many reflections from a sufficiently small unit
cell, the probability theory of random phase estimates
(Sayre equation regime) gives practical phase determination.
Certifying this formally is the remaining content of
O-CRYST.PHASE.
\end{remark}

\subsection{Updated Phase Problem Status}

\begin{definition}[\status{D}]\label{def:cryst-centrosymmetric}
\textup{(Centrosymmetric Configuration.)}
A certified configuration $X$ is \emph{centrosymmetric} when its
minimal faithful symmetry carrier $\mathcal{G}_{\mathrm{Cryst}}$
contains an inversion operation.  Equivalently: the MSC-selected
space group of $X$ belongs to one of the 73 centrosymmetric
space group types among the 230.  For a centrosymmetric
configuration, the structure factors satisfy
$F_X(hkl) \in \mathbb{R}$ (real-valued) for all $(hkl)$ in
the certified probe set --- the Friedel symmetry is exact.
\end{definition}

\begin{theorem}[\status{C}]\label{thm:cryst-centro-PHASEINVIS}
\textup{(PHASE-INVIS for Centrosymmetric Configurations.)}
\textup{[Conditional on O-CRYST.SG discharged and the certified
configuration $X$ being centrosymmetric.]}
For centrosymmetric configurations, the phase predicate takes only
the values $\phi(hkl) \in \{0, \pi\}$, and the sign of
$F_X(hkl)$ is determined by the positivity constraint
(Thm.~\ref{thm:cryst-positivity-phase}) up to a global sign
ambiguity.  This is certified structural PHASE-INVIS in the
centrosymmetric subclass: no additional probe is needed to
constrain the phases to this two-valued set, and the global sign
is not quotient-visible under intensity-only measurement.
\end{theorem}

\begin{proof}
By O-CRYST.SG, the MSC-selected space group of $X$ is determined.
For the 73 centrosymmetric space groups, the inversion symmetry
forces $F_X(hkl) = F_X(\overline{hkl})^*$, and since inversion
maps $(hkl) \to (\overline{hkl})$, both are real: $F_X(hkl) \in
\mathbb{R}$.  The phases are therefore $\phi(hkl) \in \{0, \pi\}$:
the sign of the real structure factor.  By the positivity
constraint (Thm.~\ref{thm:cryst-positivity-phase}), the admissible
sign assignments are those making
$\rho(\mathbf{r}) = \sum_{hkl} F_X(hkl)
e^{-2\pi i\mathbf{h}\cdot\mathbf{r}} \ge 0$.  This constrains
the signs but cannot eliminate the global sign ambiguity:
$\{F_X(hkl)\}$ and $\{-F_X(hkl)\}$ produce the same $\rho$ up
to reflection, giving the same intensity profile.  Therefore the
global sign is quotient-invisible under intensity-only probes:
PHASE-INVIS holds for the global sign in the centrosymmetric
subclass.
\qed
\end{proof}

\begin{definition}[\status{D}]\label{def:cryst-friedel-discipline}
\textup{(Friedel-Symmetric Intensity-Only Discipline.)}
The declared diffraction probe discipline
$\mathcal{P}_{\mathrm{diff}}$ is \emph{Friedel-symmetric} when the
certified atomic response entering the structure factors is real
(no resonant/anomalous response channel is declared), so that
$F_X(\overline{hkl}) = F_X(hkl)^*$ and hence
$|F_X(hkl)|^2 = |F_X(\overline{hkl})|^2$ for every certified
reflection.  This is the load-bearing hypothesis of the
invisibility results below: it plays the role that the
internal-factorization property plays for the LING reference
problem --- the formal declaration that the probe channel carries
no phase-sensitive response.
\end{definition}

\begin{theorem}[\status{C}]\label{thm:cryst-enantiomorph-PHASEINVIS}
\textup{(PHASE-INVIS for the Enantiomorph: Non-Centrosymmetric
Configurations.)}
\textup{[Conditional on O-CRYST.SG discharged; the certified
configuration $X$ being non-centrosymmetric
(Def.~\ref{def:cryst-centrosymmetric} failing); and
$\mathcal{P}_{\mathrm{diff}}$ being Friedel-symmetric
(Def.~\ref{def:cryst-friedel-discipline}).]}
Let $\overline{X}$ denote the inverted (enantiomorphic)
configuration.  Then $X$ and $\overline{X}$ lie in the same
observational quotient class under intensity-only probes, while
their phase profiles differ nontrivially:
$\phi_{\overline{X}}(hkl) = -\phi_X(hkl)$, and for a
non-centrosymmetric configuration this negation is not the
identity.  The enantiomorph (absolute configuration) is therefore
quotient-invisible: PHASE-INVIS holds for the conjugation
invariant in the non-centrosymmetric class.
\end{theorem}

\begin{proof}
This is the same invariance-exhibition argument as
Thm.~\ref{thm:cryst-centro-PHASEINVIS}, one level up: exhibit a
transformation preserving all probe-visible data while altering
the hidden quantity.  Inversion of the configuration sends each
certified scatterer position $\mathbf{r}$ to $-\mathbf{r}$, so
$F_{\overline{X}}(hkl) = F_X(hkl)^*$: every phase is negated,
$\phi \mapsto -\phi$.  Intensities are unchanged pointwise,
$|F_{\overline{X}}(hkl)|^2 = |F_X(hkl)|^2$, and under the
Friedel-symmetric discipline no admissible probe outcome
distinguishes the conjugate structure-factor set, since every
declared channel reads only $|F|^2$ data.  Hence $X$ and
$\overline{X}$ have identical $\mathcal{P}_{\mathrm{diff}}$
profiles.  For a non-centrosymmetric configuration the phase set
is not contained in $\{0, \pi\}$ (otherwise all structure factors
would be real, forcing the inversion symmetry excluded by
hypothesis), so $\phi \neq -\phi$ on at least one certified
reflection and the two phase profiles genuinely differ.  This is
exactly the PHASE-INVIS clause of
Def.~\ref{def:cryst-phase-predicate} for the conjugation
invariant.  In the centrosymmetric subclass the same conjugation
degenerates to the global sign already treated by
Thm.~\ref{thm:cryst-centro-PHASEINVIS}. \qed
\end{proof}

\begin{theorem}[\status{C}]\label{thm:cryst-origin-PHASEINVIS}
\textup{(PHASE-INVIS for the Origin Coset: All Configurations.)}
\textup{[Conditional on O-CRYST.LAT discharged (lattice
translation lawful in the Book~III ledger).]}
For any admissible origin shift $\mathbf{t}$, the translated
configuration $X + \mathbf{t}$ satisfies
$F_{X+\mathbf{t}}(hkl) = e^{-2\pi i\,\mathbf{h}\cdot\mathbf{t}}
F_X(hkl)$: every intensity is unchanged while every phase shifts
by the linear form $\phi(hkl) \mapsto \phi(hkl) - 2\pi\,
\mathbf{h}\cdot\mathbf{t}$.  The phase profile is therefore
quotient-invisible along the origin coset: intensity-only probes
determine phases at best modulo the linear-phase action.
\end{theorem}

\begin{proof}
Direct computation on the certified structure factors: a common
translation of all scatterers multiplies each $F_X(hkl)$ by the
unit-modulus factor $e^{-2\pi i\,\mathbf{h}\cdot\mathbf{t}}$,
which cancels in $|F|^2$.  Lawfulness of the translation as an
admissible re-description is O-CRYST.LAT.  Since the two
configurations are admissible re-descriptions of one another with
identical intensity profiles and distinct phase profiles (for
$\mathbf{t}$ not a lattice vector), the linear-phase coset is
quotient-invisible.  Combined with
Thm.~\ref{thm:cryst-enantiomorph-PHASEINVIS}, the certified
invisible group for Friedel-symmetric intensity-only probing of
non-centrosymmetric configurations contains the origin
translations together with conjugation: phases are invisible at
least up to $\phi(\mathbf{h}) \mapsto \pm\phi(\mathbf{h}) -
2\pi\,\mathbf{h}\cdot\mathbf{t}$. \qed
\end{proof}

\begin{theorem}[\status{C}]\label{thm:cryst-anomalous-PHASERECOV}
\textup{(PHASE-RECOV for the Enantiomorph under a Nominated
Anomalous Probe Extension.)}
\textup{[Conditional on: nomination of a lawful
\emph{anomalous-scattering probe extension}
$\mathcal{P}'_{\mathrm{diff}} \supset \mathcal{P}_{\mathrm{diff}}$
--- a resonant-wavelength channel declared as a certified
observable variable, sourced through $\mathbf{U}$ with no hidden
supplement (a calibration-pattern declaration, the same nomination
class as the pending diffraction standard); and the template
machinery of Thm.~\ref{thm:cryst-phase-conditional}.]}
Under $\mathcal{P}'_{\mathrm{diff}}$ the Friedel symmetry of
Def.~\ref{def:cryst-friedel-discipline} is broken: the certified
atomic response acquires an imaginary component, the Bijvoet
differences $|F_X(hkl)|^2 - |F_X(\overline{hkl})|^2$ are
generically nonzero for non-centrosymmetric $X$, and they change
sign under conjugation.  The enantiomorph is therefore
quotient-visible under $\mathcal{P}'_{\mathrm{diff}}$: PHASE-RECOV
holds for the conjugation invariant, and the MSC selects the
unique minimal faithful phase-sensitive extension among
role-equivalent lawful nominations.
\end{theorem}

\begin{proof}
With an imaginary response component the structure factors of $X$
no longer satisfy $F_X(\overline{hkl}) = F_X(hkl)^*$, so the two
members of a Friedel pair carry independent certified intensities;
their difference is an admissible probe outcome of
$\mathcal{P}'_{\mathrm{diff}}$.  Under configuration inversion
the Bijvoet difference changes sign, so $X$ and $\overline{X}$
have distinct $\mathcal{P}'_{\mathrm{diff}}$ profiles whenever at
least one certified Bijvoet difference is nonzero --- the
generic non-centrosymmetric situation, and a finitely checkable
condition on the declared reflection set.  Hence the conjugation
invariant hidden under the Friedel-symmetric discipline
(Thm.~\ref{thm:cryst-enantiomorph-PHASEINVIS}) becomes
quotient-visible.  Lawfulness of the extension is the content of
the nominated declaration and is not re-proved here; the selection
of the unique minimal faithful extension is the PHASE-RECOV clause
of the template (Thm.~\ref{thm:cryst-phase-conditional}) by the
Zorn/finite-stabilization argument.  The origin coset
(Thm.~\ref{thm:cryst-origin-PHASEINVIS}) remains invisible under
this extension, as the linear phase factor cancels in every
$|F|^2$-type channel; origin fixing is a separate declaration, not
a probe-information question. \qed
\end{proof}

\begin{openobligation}[\status{C}]\label{ob:cryst-anomalous-nomination}
\textup{(O-CRYST.PHASE-RECOV: Anomalous-Scattering Wavelength
Nominated.)}
\textup{[Conditional on Thm.~\ref{thm:cryst-anomalous-PHASERECOV}
and the declaration below.]}
A resonant-wavelength anomalous-scattering channel is hereby
nominated as the lawful phase-sensitive probe extension
$\mathcal{P}'_{\mathrm{diff}}$: an incident-radiation wavelength
declared near a certified absorption-edge of a scatterer species
present in the configuration, so that the certified atomic response
acquires its resonant (imaginary) component.  The wavelength enters
as a \emph{certified observable variable} --- a declared property of
the probe process, on the same footing as the SPEC caesium frequency
nomination (SPEC Branch, \texttt{ob:SPEC-Planck-calibration}) --- not
as an imported physical primitive.  This extension satisfies the
three lawfulness conditions: (a) it is sourced through $\mathbf{U}$
as a declared probe-response process; (b) it introduces no hidden
supplement, the resonant component being a certified response of the
already-declared scatterer species, not a new ontology; and (c) its
intensity profile (now including the Bijvoet differences) determines
the enantiomorph by Thm.~\ref{thm:cryst-anomalous-PHASERECOV}.  The
numerical wavelength value, when a specific edge is fixed, enters only
through the empirical calibration map, not through the foundational
derivation --- exactly as the caesium value does for $\hbar$.
\end{openobligation}

\begin{openobligation}[\status{C}]\label{ob:cryst-diffraction-standard}
\textup{(O-CRYST.CAL: Diffraction Standard Nominated.)}
\textup{[Conditional on the diffraction calibration pattern and the
declaration below.]}
A characteristic-emission wavelength is hereby nominated as the
diffraction length standard for the branch's metric calibration, in
the calibration-pattern slot left pending in the closure ledger:
\[
  \lambda_{\mathrm{cal}} \;=\; \lambda(\mathrm{Cu}\,K\alpha_1)
  \;\approx\; 1.5406\;\text{\AA}
  \quad(\text{declared empirical reference wavelength}).
\]
This plays for CRYST the role the caesium hyperfine frequency plays
for SPEC: a declared empirical reference that converts the certified
dimensionless diffraction-geometry ledger (Bragg-angle ratios, which
are quotient-visible without it) into SI length units.  It is sourced
through $\mathbf{U}$ as a certified emission-process observable, adds
no hidden supplement, and its numerical value lives strictly in the
empirical calibration map: the certified \emph{ratios} of lattice
spacings are nomination-independent, and only their SI \emph{scale}
is fixed by $\lambda_{\mathrm{cal}}$.  This completes the
calibration-pattern row recorded as pending for the diffraction
standard.
\end{openobligation}

\begin{openobligation}[\status{O}]\label{ob:cryst-PHASE-progress}
\textup{(O-CRYST.PHASE: Current Status and Remaining Frontier.)}
Phase problem progress summary after Sessions~L+1--L+4:

\medskip
\textbf{Conditionally established [C]:}
\begin{enumerate}[label=\textup{(\arabic*)}]
  \item Certified phase ambiguity: Babinet companion
    (Thm.~\ref{thm:cryst-babinet-invis}).
  \item Positivity constrains admissible phases
    (Thm.~\ref{thm:cryst-positivity-phase}).
  \item PHASE-INVIS for centrosymmetric configurations up to
    global sign: the 73 centrosymmetric space groups force
    real structure factors; global sign not intensity-visible
    (Thm.~\ref{thm:cryst-centro-PHASEINVIS}).
\end{enumerate}

\textbf{Conditionally established this pass [C] (continued):}
\begin{enumerate}[label=\textup{(\arabic*)}, start=4]
  \item PHASE-INVIS for the enantiomorph: conjugation negates all
    phases while preserving all Friedel-symmetric intensities;
    absolute configuration quotient-invisible
    (Thm.~\ref{thm:cryst-enantiomorph-PHASEINVIS}).
  \item PHASE-INVIS for the origin coset: linear phase factor
    cancels in every intensity; phases determined at best modulo
    the origin action (Thm.~\ref{thm:cryst-origin-PHASEINVIS}).
  \item PHASE-RECOV for the enantiomorph, conditional on the
    nominated anomalous-scattering extension: Bijvoet differences
    make conjugation quotient-visible
    (Thm.~\ref{thm:cryst-anomalous-PHASERECOV}).
\end{enumerate}

\textbf{Genuinely open [O] (narrowed):}
\begin{enumerate}[label=\textup{(\arabic*)}]
  \item \emph{PHASE-INVIS/RECOV modulo the certified invariance
    group}: for generic non-centrosymmetric configurations, decide
    whether intensities determine the phases \emph{modulo} the
    certified invisible group (origin coset together with
    conjugation), or whether further quotient-invisible phase
    freedom (homometry) persists generically.  This residual is a
    structure-factor fibration question outside the current NFC
    collar-algebra toolkit --- the same register as
    \texttt{ob:rh-s1-formal} (toolkit-boundary item).
  \item \emph{Nominations pending}: the concrete
    anomalous-scattering declaration (resonant wavelength as
    certified observable variable) and the diffraction standard
    remain calibration-pattern declarations to be made.
\end{enumerate}
\end{openobligation}

%% ---------------------------------------------------------------
\section{Referee Attack Table and Non-Jargon Architecture}
\label{sec:cryst-referee}
%% ---------------------------------------------------------------

\begin{remark}[\status{R}]\label{rem:cryst-referee-table}
\textbf{Referee attack table (Book~VII external falsifiability).}
Each step can be refuted by exhibiting the named failure:

\medskip
{\small
\noindent\begin{tabular}{|l|p{8.5cm}|}
\hline
\textbf{Step} & \textbf{Refutation condition} \\
\hline
O-CRYST.PR & probe response not sourced through $\mathbf{U}$ \\
O-CRYST.GAP & certified crystalline configuration with $B_{\mathrm{Cryst}}(X)=0$ \\
O-CRYST.LAT & translation cost not additive in Book~III ledger \\
O-CRYST.NC & fifth intensity-loss channel beyond the declared BLC package \\
O-CRYST.DOM & $\varepsilon^{\mathrm{dom}}_n \not\to 0$ \\
O-CRYST.EWALD & $\varepsilon^{\mathrm{Ewald}}_n \not\to 0$ \\
O-CRYST.SG & MSC selects a carrier outside the 230 space groups \\
O-CRYST.END & any prior discharge fails \\
O-CRYST.PHASE & certified intensity-only probe recovers the phase (PHASE-INVIS),
  or undeclared import in extended probe family (PHASE-RECOV) \\
\hline
\end{tabular}}
\end{remark}

\begin{remark}[\status{R}]\label{rem:cryst-nonjargon}
\textbf{Non-jargon proof architecture.}
The branch result in standard analytic language:

\begin{enumerate}[label=\textup{(\arabic*)}]
  \item \emph{Finite scattering approximants.}
    Exhaustion $U_1 \subset U_2 \subset \cdots$ of $\mathbb{R}^3$
    with $|\partial U_n|/|U_n| \to 0$; certified
    scattering-density variable; gauge-invariant physical
    observable $I(hkl) = |F(hkl)|^2 \ge 0$.

  \item \emph{Certified observables and admissible tests.}
    Finite probe discipline (intensity responses only, real and
    nonnegative); translation equivalence from probe data
    alone; four anti-smuggling conditions (no Euclidean
    primitive, no Fourier primitive, no atomic species, no
    chemical formula).

  \item \emph{Boundary control and no leakage.}
    Book~II unit cell boundary capacity law; Van Hove
    condition $|\partial U_n|/|U_n| \to 0$.

  \item \emph{Bragg gap (core theorem).}
    $B_{\mathrm{Cryst}}(X) > 0$ uniformly in $n$ for certified
    crystalline configurations.  Three components: (a) positivity
    of intensity at peak positions; (b) $I_{\mathrm{peak}} >
    I_{\mathrm{bg}}$ by non-trivial probe separation; (c) gap
    persistence from Book~III transport (no gap collapse without
    declared structure-change event).

  \item \emph{Space group reduction.}
    Retain only the MSC-selected minimal faithful symmetry
    carrier; remove non-load-bearing sectors (abelian factors,
    topological artifacts, hidden chemical classes).
    MSC carrier coincides with one of the 230 space groups.

  \item \emph{Domain and Fourier bridges.}
    $D_{\mathrm{Cryst}} \sim \mathbb{R}^3$ at no-loss scope;
    Bragg equation and structure factors licensed as effective
    shorthands; geometric decay
    $O(e^{-B_{\mathrm{Cryst}} \cdot n})$ for both losses.

  \item \emph{Phase problem.}
    Intensity data $\{I(hkl)\}$ is a non-injective function of
    the structure factors: Babinet pairs establish certified
    phase ambiguity.  Positivity constraint from the scattering
    density restricts the admissible phase space.  Full phase
    determination requires either a genericity argument (full
    PHASE-INVIS) or an extended probe family (PHASE-RECOV).

  \item \emph{Endpoint.}
    Certified Bragg peak structure with positive gap, additive
    lattice ledger, no-hidden-channel control, Euclidean and
    Fourier notation licensed, space group classification
    derived.
\end{enumerate}

\textbf{The two places where the NFC machinery remains
visible:} Step~4 (Bragg gap as the uniform positive lower
bound, parallel to YM uniform coercive estimate) and Step~6
(domain and Fourier no-loss bridges, parallel to YM continuum
no-loss theorem).  Everything else is standard analytic and
crystallographic language.
\end{remark}

%% ---------------------------------------------------------------
\section{Updated Closure Ledger and Obligation Register}
\label{sec:cryst-closure-updated}
%% ---------------------------------------------------------------

\begin{remark}[\status{R}]\label{rem:cryst-closure-updated}
\textbf{Post-discharge obligation register.}

\medskip
PR [C], GAP [C], LAT [C], NC [C], DOM [C], EWALD [C], SG [C], END [C],
PHASE [O].  The eight [C] obligations are discharged by
Thms.~\ref{thm:cryst-PR-discharged}--\ref{thm:cryst-END-discharged}.
O-CRYST.PHASE remains open.

\medskip
\textbf{Genuinely open obligations (3):}
\begin{enumerate}[label=\textup{(\arabic*)}]
  \item \emph{O-CRYST.PHASE (full PHASE-INVIS or PHASE-RECOV)}:
    Babinet ambiguity [C] and positivity constraint [C]
    established; full generic phase determination or certified
    extended-probe recovery not yet proved.
  \item \emph{Implicitly open}: declaration of a calibration-eligible
    diffraction standard (the crystallographic analogue of
    O-SPEC.Planck-calibration).  This would certify the declared
    probe-energy window $\mathcal{W}_{\mathrm{diff}}$ against a
    specific empirical standard.
  \item \emph{Structure determination}: even with the phase
    problem partially resolved, uniquely determining the
    three-dimensional arrangement of scatterers from the
    certified intensity and phase data requires a
    structure-determination theorem that has not yet been
    promoted to canonical form.
\end{enumerate}

\textbf{Discharge summary.} The crystallography branch has
reached conditional CERT-CLOSE on the
diffraction-periodicity-symmetry-domain endpoint.  Eight of
nine original obligations are conditionally discharged.
The phase problem is the sole remaining structural frontier.
\end{remark}



 nd{document}

\end{document}
